\documentclass[11pt]{amsart}
\IfFileExists{lmodern.sty}%
  {\usepackage[T1]{fontenc}\usepackage{lmodern}\usepackage{microtype}}%
  {\usepackage[expansion=false]{microtype}}
\usepackage{amssymb,mathtools}
\usepackage{booktabs,array}
\usepackage{enumitem}
\emergencystretch=2.5em
\usepackage[colorlinks=true,linkcolor=blue,citecolor=blue,urlcolor=blue]{hyperref}

\newcommand{\FormalizationURL}{https://github.com/slitsey/catalan-irrationality-measure}
\newcommand{\FormalizationLink}{\url{\FormalizationURL}}

% Numerical constants of the certificate.
\newcommand{\Bx}{589.19666}                             % threshold x of the criterion
\newcommand{\Rho}{301.54258}                            % perturbation exponent rho
\newcommand{\Ctwo}{4.28886}                             % cost of a fixed column, second case
\newcommand{\Cone}{4.09661}                             % cost of a fixed column, first case
\newcommand{\Utarget}{2.255858836}                      % energy bound U
\newcommand{\NormCap}{0.7736205195}                     % bound for 2|u|^2+|v|^2/2
\newcommand{\SupX}{-0.96450499}                         % sup of the row field
\newcommand{\SupY}{-1.57499400}                         % sup of the column field
\newcommand{\FieldCap}{-2.53949899}                     % SupX + SupY
\newcommand{\ResCap}{-0.0631062}                        % sup of the residual field
\newcommand{\EtaInf}{7.987762}                          % real place minus denominators (no content)
\newcommand{\EtaTot}{27.124717}                         % EtaInf + Chi
\newcommand{\Kappa}{1.956605}                           % kappa = B/rho
\newcommand{\Kcoef}{160}                                % number of Chebyshev coefficients
\newcommand{\Ngeo}{75}                                  % geometric terms
\newcommand{\Ncomplex}{41}                              % complex geometric terms
\newcommand{\Nblocks}{47, 38 and 23}                    % blocks of the three covers
\newcommand{\Ncells}{373, 302 and 176}                  % cells of the three covers
\newcommand{\Mmargin}{2.986\cdot10^{-6}}                % slack of the objective
\newcommand{\Disc}{1.645\cdot10^{-5}}                   % discriminant of the budget
\newcommand{\DiscBound}{1.656\cdot10^{-5}}
\newcommand{\TDelta}{6422.8355}                         % n_0^2 (U - log 2 / 2)
\newcommand{\EtaExact}{27.12683}                        % TDelta minus the denominator exponent
\newcommand{\Blimit}{589.19}                            % d rho / eta
\newcommand{\NCtwo}{248.75}                             % n_0 c_2
\newcommand{\PayL}{53.00012}                            % d + 4 d delta kappa
\newcommand{\PayR}{53.072}                              % eta kappa
\newcommand{\Digest}{ff307ca0287dc5224e268202cb0b09776e933652acd2f2f879492af28d3e381d}
\newcommand{\Chi}{19.136955}                            % rounded-down content area
\newcommand{\Wcap}{52.788}                              % cover constant of the weights

\newtheorem{theorem}{Theorem}[section]
\newtheorem{lemma}[theorem]{Lemma}
\newtheorem{proposition}[theorem]{Proposition}
\newtheorem{corollary}[theorem]{Corollary}
\theoremstyle{definition}
\newtheorem{definition}[theorem]{Definition}
\theoremstyle{remark}
\newtheorem{remark}[theorem]{Remark}

\numberwithin{equation}{section}

\newcommand{\Z}{\mathbb Z}
\newcommand{\Q}{\mathbb Q}
\newcommand{\R}{\mathbb R}
\newcommand{\N}{\mathbb N}
\newcommand{\F}{\mathbb F}
\newcommand{\cV}{\mathcal V}
\newcommand{\cR}{\mathcal R}
\newcommand{\cA}{\mathcal A}
\newcommand{\cL}{\mathcal L}
\newcommand{\rev}{\operatorname{rev}}
\newcommand{\lean}[1]{\texttt{#1}}

\begin{document}

\title{An irrationality measure for Catalan's constant}

\author{Sean Litsey}
\email{sean.litsey.math@gmail.com}

\date{October 24, 2026}

\subjclass[2020]{Primary 11J82; Secondary 11J72, 31A15, 68V20}

\keywords{Catalan's constant, irrationality measure, determinants, rank of reductions modulo p,
logarithmic potentials, Gauss quadrature, formal verification}

\begin{abstract}
Let $G=\sum_{j\ge0}(-1)^j(2j+1)^{-2}$ be Catalan's constant. Building on OpenAI's recent proof
that $G$ is irrational, we prove that $|G-a/b|>b^{-590}$ for every integer $a$ and every
sufficiently large integer $b$; thus the irrationality exponent of $G$ is at most $590$.
OpenAI's proof compares, for a determinant $\Delta$ whose entries are linear forms in $1$ and $G$,
a lower bound coming from the denominators of $\Delta$ with an upper bound coming from an integral
representation. The determinant depends on five integer parameters, a \emph{design}. We prove,
for an arbitrary design, that the reduction modulo $p$ of the matrix of rational parts has a
large kernel for every prime $p$ in four explicit windows, with closed formulas in the five
parameters, and that the denominator of the determinant is smaller at $p$ by the dimension of this
kernel. With the design $(58,9,3,5,3)$, a Gauss quadrature formula and an identity of Andr\'eief type
with fixed points, which bound the determinant on an exponentially small neighbourhood of $G$, and a
criterion with a free lower end of the window of primes, this gives the measure. All steps are
formalized in the Lean proof assistant, with only the standard axioms.
\end{abstract}

\maketitle

\tableofcontents

\section{Introduction}\label{sec:intro}

Catalan's constant
\[
 G=\sum_{j=0}^\infty\frac{(-1)^j}{(2j+1)^2}=L(2,\chi_{-4})=0.91596\,55941\,77219\ldots
\]
is the value at $s=2$ of the Dirichlet $L$-function of the nontrivial character modulo $4$. Its
arithmetic has long been studied by analogy with $\zeta(2)$ and $\zeta(3)$, but until recently
it was not known to be irrational. In October 2026 OpenAI \cite{OpenAI2026} proved that $G$
is irrational and verified the proof in Lean. (An earlier claimed proof \cite{Sun2026} was shown
to contain a gap \cite{Wachs2026}.)

The \emph{irrationality exponent} $\mu(\xi)$ of a real number $\xi$ is the infimum of the real
numbers $\mu$ such that $|\xi-a/b|>b^{-\mu}$ for all integers $a$ and all sufficiently large
integers $b$. Every irrational number has $\mu(\xi)\ge2$, and almost every real number has
$\mu(\xi)=2$. An \emph{irrationality measure} is an explicit upper bound for $\mu(\xi)$.

\begin{theorem}\label{thm:main}
For every integer $a$ and every sufficiently large integer $b$,
\[
 \left|G-\frac ab\right|>b^{-590}.
\]
In particular $\mu(G)\le590$.
\end{theorem}

The proof uses the prime number theorem. The theorem and every result in this paper have been
verified in the Lean~4 proof assistant on top of OpenAI's formalization, with no unproved steps and
only the standard axioms (Section~\ref{sec:formal}).

\subsection{Background and related results}
Irrationality proofs for the values of $L$-functions at integers follow, almost without exception,
the pattern of Ap\'ery's proof for $\zeta(3)$ \cite{Apery79}. One constructs integer linear forms
$q_n\xi-p_n$ that tend to zero, with $q_n\ne0$, faster than the growth of their coefficients would allow
for a rational $\xi$. If $|q_n\xi-p_n|=e^{-\alpha n+o(n)}$ and $|q_n|=e^{\beta n+o(n)}$, then
$\xi$ is irrational as soon as $\alpha>0$, and $\mu(\xi)\le1+\beta/\alpha$. Beukers
\cite{Beukers79} gave the linear forms for $\zeta(2)$ and $\zeta(3)$ as integrals of Legendre-type
polynomials. Hata \cite{Hata93} refined such integrals and introduced
the systematic use of the primes that divide all the coefficients of the forms, an idea going back
to Chudnovsky. Rhin and Viola \cite{RhinViola96,RhinViola01} found that families of integrals
for $\zeta(2)$ and $\zeta(3)$ are permuted by a finite group, and used the invariance of the integrals
to divide all their coefficients by a large common factor. This \emph{arithmetic method} gave the
sharpest measures for $\zeta(2)$ and $\zeta(3)$, and Marcovecchio \cite{Marcovecchio09} applied it to
$\log2$. For $\pi$, Zeilberger and Zudilin \cite{ZeilbergerZudilin20} obtained $\mu(\pi)\le7.1032\ldots$
from Salikhov's integrals \cite{Salikhov08} with an additional arithmetic saving.

Since September 2026 a series of new results in this area has appeared, several of them obtained with
AI systems. OpenAI proved $\mu(\pi)=2$, with a Lean formalization \cite{OpenAIPi2026}, and Liu, Jiang and
Zhang proved $\mu(\log r)=2$ for every rational $r>0$, $r\ne1$ \cite{LiuJiangZhang26}. Giraudin lowered
the bounds for $\mu(\zeta(2))$ and $\mu(\zeta(3))$ to $5.0193784$ and $5.5138800$ by applying the
Rhin--Viola method at new integral directions \cite{Giraudin26,Giraudin26z3}. Leibtag, Barkan, Kalisch,
Weinbaum, Shalyt and Kaminer \cite{LBKWSK26} showed that the classical constructions are
one-dimensional trajectories in conservative matrix fields, and by searching over other trajectories
obtained $\mu(\zeta(2))<4.62675$ and $\mu(\zeta(3))<5.4197$. Prellberg \cite{Prellberg26} proved that
$\mathrm{Li}_2(-1/4)$, $\mathrm{Li}_2(1/5)$ and $\mathrm{Li}_2(-1/3)$ are irrational with the integrals of
Rhin--Viola and Viola--Zudilin. Fauzan announced a proof that $\zeta(5)$
is irrational, using Hankel determinants \cite{Fauzan2026}, and Lai combined this method with simultaneous
approximations to show that $1$, $\zeta(3)$ and $\zeta(4)$ are linearly independent over $\Q$
\cite{Lai2026}. Several of these arguments, like the present one, use determinants of large size rather
than single sequences of linear forms; Zudilin's determinantal criterion \cite{Zudilin} is an earlier
instance of this idea.

For Catalan's constant the analogous linear forms exist but do not decay fast enough.
Rivoal and Zudilin \cite{RivoalZudilin2003} proved that at least one of the numbers
$\beta(2),\beta(4),\ldots,\beta(2m)$ is irrational for suitable $m$, where $\beta$ is the Dirichlet beta
function; $G=\beta(2)$ is not singled out. Zudilin \cite{Zudilin2003} found an Ap\'ery-like
recurrence for $G$ whose solutions give linear forms that converge to zero, but too slowly to prove
irrationality, and \cite{Zudilin02Cheb,KrattRivoal08,KrattZudilin19,Zudilin19Survey} describe
related constructions and their limits. Eskandari, Murty and Nemoto \cite{EMN2025} interpret linear forms
in Catalan's constant through mixed motives. Leibtag et al.\ \cite{LBKWSK26} found further Ap\'ery-like
approximations to $G$ in their matrix fields; numerically, the best of them still has a negative
approximation exponent, about $-0.28$.

OpenAI's argument \cite{OpenAI2026} has a different shape. Instead of a sequence of linear forms it
uses a determinant $\Delta$ of size $48N$ whose entries are rational linear combinations of $1$ and
$G$. Determinants of linear forms were used for irrationality in Zudilin's determinantal approach
\cite{Zudilin}; here the entries are double integrals against the Cauchy kernel $1/(1-ts)$. If $G$
were rational, the denominators of the entries would give a lower bound
$|\Delta|\ge e^{-(\mathsf A+o(1))n^2}$ with $n=48N$. An integral representation, of the kind that
appears in random matrix theory through Andr\'eief's identity \cite{Forrester2018Andreief}, gives the
upper bound $|\Delta|\le e^{-(\mathsf U+o(1))n^2}$, and an energy estimate shows $\mathsf U>\mathsf A$.
The margin is small: $\mathsf U-\mathsf A\approx1.3\cdot10^{-4}$. The nonvanishing of $\Delta$ along the
primes $N=p$ is proved separately by reduction modulo $p$.

The first irrationality measure for $G$ was obtained by Jha \cite{Jha2026}, who proved
$\mu(G)\le10^8$. Carreira \cite{Carreira2026} treated the five integer parameters of the construction
as free, stated conditions under which OpenAI's argument goes through for general parameters, and
obtained $\mu(G)<2294$ in the revised version of his note; he also formulated the best real-place
constant of the method as an equilibrium problem for two measures, and in \cite{Carreira2026b} he
records attempts to carry the method to $\beta(4)$. His bound is the closest to ours, and the
comparison is instructive. The determinants are of the same kind. Carreira makes the determinant
small near $G$ by bounding all its minors, as Jha does; here a quadrature formula for the coefficient of
the constant is used instead (Section~\ref{sec:real}). And the present paper uses, for every prime $p$
in an interval of length about $14N$, the rank of the reduction of the matrix modulo $p$; this lowers
the denominator of the determinant by more than twice the margin that remains between the real and the
finite places without it (Section~\ref{sec:intro-ideas}). Eskandari \cite{Eskandari2026} studies rational
approximations to $G$ from another point of view. For $L(2,\chi_{-3})$, Calegari, Dimitrov and Tang
proved linear independence of $1$, $\zeta(2)$ and $L(2,\chi_{-3})$ \cite{CDT2024}, and their
holonomy bounds give $\mu(L(2,\chi_{-3}))\le24781$ \cite[Thm.~5.1]{CDT2025}. A companion paper
\cite{Litsey2026b} proves measures for $G\pm5\pi^2/48$ and $L(2,\chi_{-3})$ with a related method.

\subsection{What a measure requires}\label{sec:intro-need}
Write $\Delta_N(X)$ for the determinant regarded as a polynomial in the constant, so that
$\Delta_N(G)$ is OpenAI's determinant. Three properties of the integer polynomials $Q_N$ obtained
from $\Delta_N$ by clearing denominators turn the comparison of OpenAI into a measure.
\begin{enumerate}[label=(\roman*)]
\item \emph{Smallness on a neighbourhood of $G$.} The upper bound must hold not only at $X=G$, but for
all $X$ with $|X-G|\le e^{-\rho N}$, and the exponent $\rho$ enters the final bound linearly.
\item \emph{Uniform nonvanishing.} $Q_p(a/b)\ne0$ for primes $p\nmid b$ (here $p^2\nmid b$), in place
of $Q_p(G)\ne0$.
\item \emph{Arithmetic.} The integer $b^{\deg Q_N}Q_N(a/b)$ is nonzero, hence at least $1$ in absolute
value; the smaller the denominator of $\Delta_N$ and the smaller its degree, the stronger the
conclusion. The prime $p$ must be chosen not dividing $b$, and a naive choice costs a factor.
\end{enumerate}
With the degree $d$ of $Q_N$ in $X$, the decay rate $\eta$ of $Q_N$ (normalised by $N^2$) and $\rho$, the
conclusion is $\mu(G)\le d\rho/\eta$ (Proposition~\ref{prop:criterion}). Everything below serves to make
$d$ and $\rho$ small and $\eta$ large for a design.

\subsection{New ideas}\label{sec:intro-ideas}
\begin{enumerate}[label=\arabic*.,wide]
\item \emph{All designs at once.} The determinant has five integer parameters
$(n_0,b_0,q_0,g_0,h_0)$, a \emph{design}, and every argument below is carried out for an arbitrary
design, with the parameters as symbols: the degree $n-g$, the denominators at all primes, the rank
deficiencies, the nonvanishing and the real-place budgets. A design is then data. We use
$(58,9,3,5,3)$, so that $n=58N$, $b=9N$, $q=3N$, $g=5N$, $h=3N$ and the degree is $53N$.

\item \emph{Rank of the reduction modulo $p$, in closed form.} Let $\nu_N(p)$ be the dimension of the
kernel (on the left) of the reduction modulo $p$ of $p^2$ times the matrix of rational parts. We prove
four lower bounds for $\nu_N(p)$ (Theorem~\ref{thm:content}), each a piecewise linear function of $p$
with coefficients that are linear in the five parameters, each on its own explicit window of primes:
antisymmetric combinations of rows centred at $g$, $g+p$ and $g+2p$ (Theorem~\ref{thm:T});
a reflection of Chebyshev polynomials modulo $p$ (Theorem~\ref{thm:U}); and two block arguments on
columns, using Frobenius relations $(1-X)^p=1-X^p$ for the coefficients of the row polynomials
(Theorems~\ref{thm:U6} and~\ref{thm:S}). They are proved from one digit formula for the entries
modulo $p$ (Lemma~\ref{lem:digits}). For $(58,9,3,5,3)$ they cover the primes $12N<p\le26N$.

\item \emph{One bound per prime.} At a prime $p$ the valuation of $\Delta_N(z)$ is at least
$\nu_N(p)-2n$ for every $z$ that is $p$-integral (Proposition~\ref{prop:valuation}); the denominator
of the determinant is therefore divided by $p^{\nu_N(p)}$, and the best of the four bounds is used
once. No separate common factor of the coefficients is needed: integer-valuedness of $\Delta_N$ at
the integers already makes the coefficients of $Q_N$ divisible by what is gained
(Proposition~\ref{prop:denominators}). The gain is $\int\kappa(x)\,dx=35369/1848$ in units of $N^2$.

\item \emph{Perturbation by quadrature with an odd number of nodes.} The smallness on a neighbourhood
of $G$ is obtained from a Gauss-type quadrature formula for the coefficient of the constant, with
$d_0N=53N$ nodes (an odd number for odd $N$) kept at a distance from the origin, and an identity of
Andr\'eief type with fixed points (Section~\ref{sec:real}).

\item \emph{Trial potentials with complex poles.} The energy estimate at the real place uses trial
sequences whose generating functions are logarithmic potentials of point charges at complex points
close to the edges of the support of the extremal measure, and a near-field estimate at the first
Chebyshev coefficient (Section~\ref{sec:real}).

\item \emph{A criterion with a free lower end of the window.} In the choice of the prime, the window
of candidates is $(V,U]$ with $U=\kappa\log b$ and $V=(\kappa-\omega)\log b$, and the natural choice
$\omega=1$ requires $\kappa>2$. We prove the criterion for every $\omega\in(0,1]$ with the hypothesis
$\kappa>2\omega$ (Proposition~\ref{prop:criterion}); the design above has $\kappa=\Kappa$ and uses
$\omega=3/4$.
\end{enumerate}

\subsection{The constants}
For the design $(58,9,3,5,3)$, in the unit $N$ with $n=58N$: the degree of $Q_N$ is $d=53N$; the
perturbation exponent is $\rho=\Rho$; the decay rate of $Q_N$ is $\eta=\EtaTot$, of which the real
place minus the denominators without the gain is $\EtaInf$ and the gain from the rank deficiencies is
$\Chi$ (rounded down from $35369/1848=19.13906\ldots$). Decimals are truncated; the constants of the
formal proof are the exact rational numbers. The criterion gives the exponent
$\Bx$ as the least $B$ allowed by its conditions
(Section~\ref{sec:proof}); the integer $590$ is the exponent of Theorem~\ref{thm:main}, for which
all the conditions are verified in Lean with rational constants.

\subsection{Organization and notation}
Section~\ref{sec:overview} sets up designs and states the rank deficiencies, the valuations, the
denominators and the criterion, which are the new ideas in precise form. Section~\ref{sec:setup}
recalls OpenAI's construction for a design and proves the degree, the denominators and the
nonvanishing. Section~\ref{sec:content} proves the rank deficiencies. Section~\ref{sec:real} contains
the estimates at the real place, Section~\ref{sec:proof} proves Theorem~\ref{thm:main}, and
Section~\ref{sec:remarks} makes remarks on the method. Section~\ref{sec:formal} describes the
formalization; the disclosure of AI use follows it.

We use the notation of OpenAI's paper \cite{OpenAI2026} for the construction and refer to it for
details; citations of the form \cite[Prop.~2.1]{OpenAI2026} point to its numbered statements. We
write $\log$ for the natural logarithm, $v_p$ for the $p$-adic valuation, $\theta(x)=\sum_{p\le x}\log p$
for Chebyshev's function, $T_k$ and $U_k$ for the Chebyshev polynomials of the first and second kind
(with $U_{-1}=0$ and $T_{-k}=T_k$, $U_{-k-1}=-U_{k-1}$), $[X^k]F$ or $F[k]$ for the coefficient of $X^k$
in a polynomial $F$, $y_+=\max(y,0)$, and $\operatorname{rev}_E F(X)=X^EF(1/X)$ for the reversal of a
polynomial of degree at most $E$.

\section{Designs, rank deficiencies and the exponent}\label{sec:overview}

\subsection{Designs}
A \emph{design} is a quintuple $(n_0,b_0,q_0,g_0,h_0)$ of integers with $n_0,b_0,q_0,h_0\ge1$,
$g_0\ge0$ and $2g_0<n_0$. For $N\ge1$ we write
\begin{gather*}
 n=n_0N,\quad b=b_0N,\quad q=q_0N,\quad g=g_0N,\quad h=h_0N,\\
 L=n+b+q,\quad C=L+g,\quad H=C+h,\quad A=b+q+2g ,
\end{gather*}
and $d_0=n_0-g_0$. OpenAI's choice is the design $(48,7,4,4,2)$. We work throughout with
\begin{gather*}
 D=(58,9,3,5,3):\quad n=58N,\quad b=9N,\quad q=3N,\quad g=5N,\quad h=3N,\\
 L=70N,\quad C=75N,\quad H=78N,\quad A=22N,
\end{gather*}
and $d_0=53$. The construction of OpenAI's determinant for a design is recalled in
Section~\ref{sec:setup}; it yields an $n\times n$ matrix $\cA_N=\cA^0_N+G\cL_N$ with rational matrices
$\cA^0_N$ and $\cL_N$, and the polynomial $\Delta_N(X)=\det(\cA^0_N+X\cL_N)$ of degree at most $n-g$
(Lemma~\ref{lem:degree}). The search that led to $D$ among designs, scored by the exponent of
Proposition~\ref{prop:criterion}, is not part of the proof.

\subsection{Rank deficiencies}
For an odd prime $p$ with $H<p^2$ the entries of $p^2\cA^0_N$ are $p$-integral (Lemma~\ref{lem:digits});
let $\cA^0_{N,p}$ be its reduction modulo $p$, and let $\nu_N(p)$ be the dimension of its left kernel, that
is $n$ minus its rank. Write $\Lambda=g+q+h$ and $e=(p+1)/2-q-h$.

\begin{theorem}\label{thm:content}
Let $D$ be a design, $N\ge1$, and $p$ an odd prime with $H<p^2$. Then $\nu_N(p)$ is at least each of the
following quantities, whenever the stated conditions hold; negative values are replaced by $0$.
\begin{enumerate}[label=(\alph*)]
\item (Theorem~\ref{thm:T}) If $2g<p$, $H\le4p$ and $D_T=\min(g,\,n-1-g-2p,\,C-3p)\ge1$:
\[
 D_T-(n+q+g+h-3p)_+ .
\]
\item (Theorem~\ref{thm:U}) If $2h\le p+1$, $L\le3p\le H$, $n+q+g+h\le3p$ and $2p\le n+q$:
\[
 \min(n,\,2p+2g+1)-g-2p-1 .
\]
\item (Theorem~\ref{thm:U6}) If $e\ge0$, $n\le4p+2g+1$, $L\le6p$ and $b+q\le p+\mu_0$, with
$\mu_0=\max(q,H+q-6p)$:
\[
 \min\Bigl(\frac{p-1}2-g,\ n+b+q+h-\frac{11p+1}2\Bigr).
\]
\item (Theorem~\ref{thm:S}) If $e\ge0$ and $n\le4p+2g+1$: $m_1-m_0-e$, with
\[
 m_0=\max(b+q,\,L+\Lambda-4p),\qquad m_1=\min(p,\,L-3p,\,b+q+2g+1).
\]
\end{enumerate}
\end{theorem}

These are the statements of the Lean theorems \lean{t\_rank}, \lean{u\_rank}, \lean{u6\_rank\_closed} and
\lean{s\_rank} (the last two in the form $m_1-m_0-e$ and its closed form), proved for every
admissible design. For $D$ they combine, for $12N<p\le26N$ and up to the loss of one unit at the
integer rounding, into the following function $\kappa$ of $x=p/N$ (units of $N$); on each cell the
value is the largest of the bounds that apply.
\begin{center}
\begin{tabular}{@{}lll@{}}
\toprule
interval of $x$ & $\kappa(x)$ & bound \\
\midrule
$(12,13]$ & $x/2-5$ & (c), $(p-1)/2-g$ \\
$(13,\,146/11]$ & $73-11x/2$ & (c), $n+b+q+h-(11p+1)/2$ \\
$(50/3,\,69/4]$ & $9x/2-75$ & (d), $(9p-2C)/2$ \\
$(69/4,\,35/2]$ & $x/2-6$ & (d), $(p-1)/2+h-b$ \\
$(35/2,\,128/7]$ & $64-7x/2$ & (d), $(2(n+q+h)-7p)/2$ \\
$(64/3,\,23]$ & $3x-64$ & (a), $3p-(n+q+h)$ \\
$(23,\,24]$ & $5$ & (a) up to $x=70/3$, then (b); in both the cap $g$ \\
$(24,\,26]$ & $53-2x$ & (b), $n-g-2p$ \\
\bottomrule
\end{tabular}
\end{center}
and $\kappa=0$ elsewhere. The area of $\kappa$ is
\[
 \int_0^\infty\kappa(x)\,dx=\frac54+\frac9{44}+\frac{49}{64}+\frac{43}{64}+\frac{121}{112}+\frac{25}6+5+6
 =\frac{35369}{1848}=19.13906\ldots
\]
The families (a)--(d) come from different arguments, and the kernels they produce are not shown to be
independent; at a prime the best bound is used once, never the sum.

\subsection{Valuations, denominators and the exponent}
The rank deficiencies enter through the following statement, which is the finite-place input.

\begin{proposition}\label{prop:valuation}
Let $p$ be an odd prime with $H<p^2$ and let $z\in\Q$ have denominator prime to $p$. Then
$p^2(\cA^0_N+z\cL_N)$ has $p$-integral entries, its reduction modulo $p$ equals $\cA^0_{N,p}$, and
\[
 v_p\bigl(\Delta_N(z)\bigr)\ge\nu_N(p)-2n .
\]
\end{proposition}

At the other primes there are the valuation bounds of OpenAI's method (Section~\ref{sec:setup}), and
at $2$ a bound from the expansion into minors. They combine to the denominators.

\begin{proposition}\label{prop:denominators}
For the design $D$ there are positive integers $D_N$ such that $D_N\Delta_N(i)\in\Z$ for every integer
$i$, such that $q^{2n}\mid D_N$ for every prime $N<q\le2N$, and such that for every $\varepsilon>0$ and
all large $N$
\[
 \log D_N\le(\mathsf A+\varepsilon)N^2,\qquad
 \mathsf A=\frac{2673}8\log2+116+\frac{24269}4-\frac{35369}{1848}=6395.7087\ldots
\]
\end{proposition}

Here $2673/8\,\log2$ is the contribution of the prime $2$, and $116+24269/4$ is the area of the
piecewise linear function that bounds the poles at the odd primes (Section~\ref{sec:setup}); the last
term is the gain $\int\kappa$ of Theorem~\ref{thm:content}.

At the real place we prove the following in Section~\ref{sec:real}. Put
\begin{equation}\label{eq:rho-def}
 \rho=n_0c_2+w,\qquad \rho=\Rho\ \text{ for }D,
\end{equation}
where $c_2=\Ctwo$ is the cost of a fixed column in the energy estimate and $w=52.7887$ is the logarithmic
size of the weights of the quadrature formula.

\begin{theorem}\label{thm:real}
Let $U=\Utarget$. For all large $N$ and all real $x$ with $|x-G|\le e^{-\rho N}$,
\[
 |\Delta_N(x)|\le\exp\bigl(-(n_0^2(U-\tfrac12\log2)-\varepsilon)N^2\bigr),\qquad n_0=58 .
\]
\end{theorem}

Since $n_0^2(U-\frac12\log2)=\TDelta\ldots$ and $\mathsf A=6395.7087\ldots$, the polynomial
$Q_N=(d_0N)!\,D_N\Delta_N$ has integer coefficients, degree at most $d_0N$, and
\[
 |Q_N(x)|\le e^{-\eta N^2}\qquad(|x-G|\le e^{-\rho N}),\qquad \eta=\EtaTot ,
\]
where $\eta$ is a rational number slightly below $n_0^2(U-\frac12\log2)-\mathsf A=\EtaExact\ldots$ The two
remaining properties of Section~\ref{sec:intro-need} are the nonvanishing of $Q_p(a/b)$ for $p^2\nmid b$
(Proposition~\ref{prop:nonvanishing}) and the divisibility of the $i$-th coefficient of $Q_N$ by $q^{2i}$
for $N<q\le2N$ (Lemma~\ref{lem:recycling}), which come from the denominators and are what is used to
choose the prime. The criterion is the following.

\begin{proposition}\label{prop:criterion}
Let $\xi\in\R$, $d\in\N$, and let $\eta,\rho,\kappa,\delta,\omega>0$ satisfy
\[
 \omega\le1,\quad 2\omega<\kappa,\quad \eta<2d,\quad 1<2\bigl(\omega-\delta(2\kappa-\omega)\bigr),\quad
 d+4d\delta\kappa<\eta\kappa .
\]
Let $Q_p\in\Z[X]$ ($p$ prime) satisfy, for all large $p$: $\deg Q_p\le dp$; $Q_p(a/b)\ne0$ whenever
$a\in\Z$, $b\ge1$, $p^2\nmid b$; $q^{2i}$ divides the $i$-th coefficient of $Q_p$ for all primes
$p<q\le2p$ and $i\le dp$; and $|Q_p(x)|\le e^{-\eta p^2}$ for $|x-\xi|\le e^{-\rho p}$. Then
\[
 |\xi-a/b|>b^{-\rho\kappa}
\]
for all integers $a$ and all sufficiently large $b$.
\end{proposition}

The conditions say, in the order of the proof of Section~\ref{sec:proof}, that the window of primes
contains a prime not dividing $b$, that the primes of the window above the chosen one are recycled, and
that the final inequality is contradictory. For the design $D$ one takes $d=53$, $\omega=3/4$,
$\delta=2.9\cdot10^{-7}$ and $\kappa=B/\rho$ with $B=590$; the last condition then reads
$B>x:=d\rho/(\eta-4d\delta)=\Bx$.

\subsection{A reading of the exponent}
The criterion shows that, for a fixed design, the bound is $d\rho/\eta$ in the limit, and that its three
ingredients come from three different places: the degree $d=n-g$ from the construction, the decay rate
$\eta$ from the real place, the denominators and the rank deficiencies, and the perturbation exponent
$\rho$ from the quadrature. The rank deficiencies matter because $\eta$ is a small difference of two
large numbers, $\TDelta$ and $6395.71$; the gain $19.14$ is more than two thirds of $\eta=\EtaTot$.

\section{OpenAI's construction for a design}\label{sec:setup}

We recall the construction of \cite{OpenAI2026} in the generality in which we use it, and we prove for an
arbitrary design the degree bound, the bounds for the denominators at the finite places and the
nonvanishing at the primes $N=p$. Carreira \cite{Carreira2026} checks that the argument of OpenAI goes
through for designs satisfying a few conditions; we need the same for the finite places and, in addition,
for the entries modulo $p^2$.

\subsection{The determinant}
Let $D=(n_0,b_0,q_0,g_0,h_0)$ be a design, with the notation of Section~\ref{sec:overview}. For $0\le r<n$
the row polynomials are
\begin{equation}\label{eq:rows}
 P_r(t)=(1-t)^ht^{C-1}T_{|r-g|}(1/t),\qquad \deg P_r<H,
\end{equation}
together with companion polynomials $D_r$ built from $U_{|r-g|-1}$ \cite[(3)]{OpenAI2026}. Thus
$P_r$ depends only on $|r-g|$. Put $f(t)=\sqrt{1-t^2}$ and $d\mu(t)=|t|\,dt/f(t)$ on $(-1,1)$, and
\[
 A_r(t)=P_r(t)-\frac32\,\mathbf 1_{\{t>0\}}\frac{f(t)}{t}D_r(t),\qquad
 \psi_k(s)=s^{b+k}(1-s)^q .
\]
The determinant studied by OpenAI is that of the $n\times n$ matrix
\begin{equation}\label{eq:matrix}
 \cA_N(r,k)=\int_{-1}^1\int_0^1 A_r(t)\,\frac{\psi_k(s)}{1-ts}\,ds\,d\mu(t)
 \qquad(0\le r,k<n).
\end{equation}
The factor $(1-s)^q$ in $\psi_k$ amounts to a finite difference of order $q$ applied to the columns.
The entries are rational linear combinations of $1$ and $G$: there are rational matrices $\cA^0_N$
and $\cL_N$ with
\begin{equation}\label{eq:L}
 \cA_N=\cA^0_N+G\,\cL_N,\qquad
 \cL_N(r,k)=\sum_{v=0}^{q}(-1)^v\binom qv\sum_{i}[t^i]P_r\cdot4\,c_{(i-b-k-v)/2},
\end{equation}
where $c_l=4^{-l}\binom{2l}{l}$ for integers $l\ge0$ and $c_z=0$ otherwise
\cite[Prop.~2.1 and (20)]{OpenAI2026}. The rational part is
\begin{equation}\label{eq:rational-part}
 \begin{gathered}
 \cA^0_N(r,k)=\sum_{v=0}^q(-1)^v\binom qv R(r,b+k+v),\\
 R(r,j)=\sum_{i<H}P_r[i]\,M(i,j)-\frac32\sum_{i<H}D_r[i]\,Z(i,j)
 \end{gathered}
\end{equation}
with rational numbers $M(i,j)$ (moments of $\mu$ against the Cauchy kernel) and $Z(i,j)$, where
$Z(i,i)=-\sum_{k\le i}k^{-2}$ and $Z(i,j)=(H_i-H_j)/(i-j)$ for $i\ne j$, $H_i=\sum_{k\le i}1/k$
\cite[Prop.~2.1]{OpenAI2026}. The design enters \eqref{eq:rational-part} only through the integer
coefficients $P_r[i],D_r[i]$ and the range of $j$, which is $[b,L)$. Put
$\Delta_N(X)=\det(\cA^0_N+X\cL_N)\in\Q[X]$, so that $\Delta_N(G)=\det\cA_N$.

\subsection{The degree}
\begin{lemma}\label{lem:degree}
$\deg\Delta_N\le n-g$.
\end{lemma}

\begin{proof}
By \eqref{eq:L} the row $r$ of $\cL_N$ depends only on $P_r$, and hence only on $|r-g|$. So the rows
$r$ and $2g-r$ of $\cL_N$ agree for $0\le r<g$, and $2g-r\le n-1$ because $2g<n$. Subtracting row $r$ from
row $2g-r$ for each $r<g$ does not change the determinant and makes $g$ rows of the new coefficient of
$X$ vanish. A determinant $\det(B+XL')$ in which $L'$ has $n-g$ nonzero rows has degree at most $n-g$ in $X$.
\end{proof}

\subsection{Rank deficiency and valuation}
We use repeatedly the following elementary fact.

\begin{lemma}\label{lem:rank-valuation}
Let $M$ be an $n\times n$ matrix with entries in $\Z_{(p)}$ whose reduction modulo $p$ has rank at most
$n-e$. Then $v_p(\det M)\ge e$.
\end{lemma}

\begin{proof}
Factor the reduction through $\F_p^{n-e}$ and lift the two factors to integral matrices $U,V$, so that
$M=UV+pE$ with $E$ integral at $p$. Expanding the determinant multilinearly in the columns, every term
that uses more than $n-e$ columns of $UV$ vanishes, and every other term has at least $e$ factors $p$.
\end{proof}

\subsection{The digit formula}
The reduction of $p^2\cA^0_N$ modulo $p$ is described by the digits in base $p$ of the coefficients of the
row polynomials. This is the one place where the arithmetic of the moments enters; it is design-free.
For an odd prime $p$ and $c\in\F_p^n$ put
\[
 P_c=\sum_rc_rP_r,\qquad D_c=\sum_rc_rD_r,\qquad W=(1-X^2)^{(p-1)/2},
\]
(reading the $P_r,D_r$ as polynomials in $X$), and let $F[k]$ denote the coefficient of $X^k$.

\begin{lemma}[digit formula]\label{lem:digits}
Let $p$ be an odd prime with $H<p^2$, let $0\le j<H$, and write $j=Jp+\ell$ with $0\le\ell<p$. Then
$p^2R(r,j)$ is $p$-integral for every row $r$, and for every $c\in\Z^n$
\[
 p^2\sum_rc_rR(r,j)\equiv\sum_{u\ge0}\Bigl(\alpha_{u,J}\,(XWP_c)[up+\ell]+\beta_{u,J}\,D_c[up+\ell]\Bigr)
 \pmod p ,
\]
where $\alpha_{u,J}=M(u-1,J)$ and $\beta_{u,J}=-\tfrac32Z(u,J)$ are $p$-integral numbers that do not depend on
$c$, on $\ell$ or on the design.
\end{lemma}

\begin{proof}
Two facts about the numbers $M$ and $Z$ are needed, for $0\le i,j<p^2$. Write $j=Jp+\ell$ and let
$\varrho\in[0,p)$ be the residue of $j-i-1$ modulo $p$, $u=(i+1+\varrho-\ell)/p$. First, $p^2M(i,j)$ is
$p$-integral and $p^2M(i,j)\equiv W[\varrho]\,M(u-1,J)\pmod p$; this is the digit reduction of the moments in
\cite[\S3]{OpenAI2026}, obtained from the $p$-adic analysis of their recursions, and it does not involve a design.
Second, $p^2Z(i,j)$ is $p$-integral and is congruent to $Z(\lfloor i/p\rfloor,J)$ if $i\equiv j\pmod p$ and to
$0$ otherwise. For the second, write $i'=\lfloor i/p\rfloor$. Since $i<p^2$, $H_i=p^{-1}H_{i'}+S_i$ and
$H^{(2)}_i=p^{-2}H^{(2)}_{i'}+S'_i$ with $S_i,S'_i$ $p$-integral, which gives the case $i=j$. If $i\not\equiv j$,
then $p^2Z(i,j)=p\cdot pZ(i,j)\equiv0$. If $i\equiv j$, $i\ne j$, then $i-j=(i'-J)p$, and the integers
$k\in(j,i]$ prime to $p$ form $i'-J$ complete systems of residues modulo $p$, so that
$\sum_k1/k\equiv(i'-J)\sum_{r=1}^{p-1}1/r\equiv0\pmod p$; hence $H_i-H_j\equiv(H_{i'}-H_J)/p$ and
$p^2Z(i,j)\equiv Z(i',J)$.

Now use \eqref{eq:rational-part}. By the first fact the part $\sum_iF[i]\,p^2M(i,j)$ with $F=P_c$ is
$\sum_iF[i]\,W[\varrho(i)]\,M(u(i)-1,J)$. For fixed $u$ the indices with $u(i)=u$ are those for which
$\varrho=up+\ell-1-i$ lies in $[0,p-1]$, and $W$ has degree $p-1$, so the sum over these $i$ is
$(XFW)[up+\ell]$; this gives the $\alpha$-terms. By the second fact the part with $D_c$ is
$\sum_{i\equiv\ell}D_c[i]\,Z(\lfloor i/p\rfloor,J)=\sum_uZ(u,J)D_c[up+\ell]$, which gives the $\beta$-terms.
The $p$-integrality of $M(u-1,J)$ and $Z(u,J)$ for $u,J<p$ is part of the first and second facts.
\end{proof}

For the filtered columns we use the following form. Let $0\le k<n$, $s=b+k$, $J=\lfloor s/p\rfloor$ and
$\mu=(s\bmod p)+q$. If $\mu<p$, the $q+1$ raw columns $s,\dots,s+q$ that enter column $k$ have the same
digit $J$ and residues $(s\bmod p)+v<p$. Since
$\sum_{v=0}^q(-1)^v\binom qvF[t+v]=[X^{t+q}]\bigl((X-1)^qF\bigr)$ for every polynomial $F$,
Lemma~\ref{lem:digits} gives
\begin{equation}\label{eq:filtered-digits}
 p^2\sum_rc_r\cA^0_N(r,k)\equiv\sum_{u\ge0}\Bigl(\alpha_{u,J}\,Q_c[up+\mu]+\beta_{u,J}\,\widehat D_c[up+\mu]\Bigr),
\end{equation}
where $Q_c=(X-1)^q\,XWP_c$ and $\widehat D_c=(X-1)^qD_c$. A combination $c$ whose responses vanish on all
raw columns in $[b,L)$ lies in the left kernel of $\cA^0_{N,p}$, and a combination whose responses on the
columns outside a set vanish and which kills the set by \eqref{eq:filtered-digits} does as well.

\subsection{Integrality and the valuation at $p$}
\begin{proof}[Proof of Proposition~\ref{prop:valuation}]
By Lemma~\ref{lem:digits} the entries of $p^2\cA^0_N$ are $p$-integral, since the filter is an integer
combination of raw columns with $j<L\le H$. The entries of $\cL_N$ are $p$-integral for odd $p$, since
$P_r[i]\in\Z$ and $c_l\in\Z[1/2]$, so $p^2\cL_N\equiv0\pmod{p^2}$ and, for $z$ with denominator prime to
$p$, $p^2(\cA^0_N+z\cL_N)$ has $p$-integral entries with the same reduction as $p^2\cA^0_N$. By
Lemma~\ref{lem:rank-valuation} the determinant of the latter has valuation at least $\nu_N(p)$, and
$\det\bigl(p^2(\cA^0_N+z\cL_N)\bigr)=p^{2n}\Delta_N(z)$.
\end{proof}

\subsection{The denominators}
Let $p$ be an odd prime with $H<p^2$ and $p\le H$. The method of \cite[\S3]{OpenAI2026}, applied to a
design, bounds from below the valuation of $\Delta_N(z)$ for every $z$ with denominator prime to $p$ by
$-\ell_N(p)$, where $\ell_N(p)$ counts the columns whose entries have a pole at $p$, once the coefficients of
the row polynomials are grouped by their digits in base $p$ and the column filter is taken into account. It
is given by two formulas, in which all differences are truncated at $0$. For $2p\le H$,
\[
 \ell_N(p)=2n-\Bigl(\min\bigl(p,\,L-p,\,A\bigr)-\max\bigl(b,\,H-2p\bigr)-q\Bigr),
\]
and for $H<2p\le2H$, with $R=(L-b)-(\min(A,L-p)-b)-(\min(p,L)-\max(b,H-p))$,
$P'=\min(A,L-p)-b$ and $E'=\min(b,H-p)-(L-p)$,
\[
 \ell_N(p)=\min\bigl(2n,\ n+R,\ 2R+P'+E'\bigr).
\]
These are the exponents of the generic valuation theorems in the formalization; for the primes that
interest us they are the same as in OpenAI's paper for its design. For the design $D$ the function
$\ell_N(p)=Nf(p/N)$ is piecewise linear in $x=p/N$ ($\lfloor\cdot\rfloor$ rounding aside):
\[
 f(x)=\begin{cases}
 116&x\le59/2,\\
 175-2x&59/2\le x\le69/2,\\
 106&69/2\le x\le39,\\
 184-2x&39\le x\le81/2,\\
 265-4x&81/2\le x\le48,\\
 217-3x&48\le x\le69,\\
 148-2x&69\le x\le70,\\
 78-x&70\le x\le78,
 \end{cases}
\]
and $f=0$ for $x\ge78$, since for $p>H$ the matrix $\cA^0_N+z\cL_N$ is $p$-integral. Its integral over
$[1,78]$ is $24269/4$, and $f=116=2\cdot58$ on $[0,1]$. Note that $f(x)=2n_0$ for $x\le59/2$, which
contains all the primes in $(N,2N]$ and, in particular, the whole range $[12,26]$ of the rank
deficiencies of Theorem~\ref{thm:content}; there the bound of Proposition~\ref{prop:valuation} improves
$-\ell_N(p)=-2n$ to $\nu_N(p)-2n$.

The primes $p\le\sqrt H$ are not covered by these bounds. For them we use the cruder bound
$6n\log_pH$ for the exponent, which contributes $O(\sqrt N\,N\log N)=o(N^2)$ in total.

\subsection{The prime $2$}
At $p=2$ the entries are not integral and the bound comes from the expansion of the determinant into
minors, in the manner of \cite[\S3]{OpenAI2026}. Its exponent is governed by a quadratic form in the
numbers $m,\ell$ of columns of two exceptional types, which is
\[
 \tfrac32m^2+\tfrac12\ell^2-(H+C-b)\,m-C\,\ell+Cn ,
\]
and the following identity shows that it is at least $-\mathsf A_2$ whenever $m,\ell\ge0$,
$m+\ell\le n$ and $C\ge n$, where
\begin{equation}\label{eq:A2}
 \mathsf A_2=\frac{(n+H-b)^2}8-\frac{n^2}2 .
\end{equation}
With $s=n-m-l$ and $m_*=(n+H-b)/4$,
\[
 \tfrac32m^2+\tfrac12\ell^2-(H+C-b)m-C\ell+Cn+\mathsf A_2
 =2(m-m_*)^2+s\,(C-n+m)+\tfrac12s^2\ge0 .
\]
This gives a $2$-adic bound $v_2(\Delta_N(i))\ge-\mathsf A_2-O(N\log N)$, uniform in the integer $i$. For $D$,
$\mathsf A_2=(58+78-9)^2/8-58^2/2=2673/8$ in units of $N^2$.

\subsection{The common denominator}
\begin{proof}[Proof of Proposition~\ref{prop:denominators}]
Let $\kappa_N(p)$ be the largest of the bounds of Theorem~\ref{thm:content} that apply at $p$ (and $0$ if
none does), for primes $p\le H$ with $H<p^2$, and put $e_N(p)=\min\bigl(\ell_N(p),\,2n-\kappa_N(p)\bigr)$. For the
primes with $p^2\le H$ use the crude bound above, for $p=2$ the $2$-adic bound, and let $D_N$ be the
product of $p^{e_N(p)}$ over $p\le H$. By Proposition~\ref{prop:valuation} and the bounds of this section,
$v_p(\Delta_N(i))\ge-e_N(p)$ for every integer $i\ge0$ and every $p$ (for $p>H$ the matrix is
$p$-integral), so $D_N\Delta_N(i)\in\Z$. A prime $q\in(N,2N]$ has $\kappa_N(q)=0$, since no window
contains it, and $\ell_N(q)=2n$, so $q^{2n}\mid D_N$ (here $N\ge79$ guarantees $q>\sqrt H$).

For the size of $D_N$, the sum of $\ell_N(p)\log p$ over the odd primes is, by the prime number theorem in
the weighted form
\begin{equation}\label{eq:weighted-pnt}
 \frac1N\sum_{p}F(p/N)\log p\longrightarrow\int_0^\infty F(x)\,dx\qquad(N\to\infty),
\end{equation}
for bounded, Riemann integrable, compactly supported $F$, equal to $N^2\int_0^{78}f(x)\,dx+o(N^2)$ with
$\int_0^{78}f=116+24269/4$. The primes with $\kappa_N(p)>0$ lie in $(12N,26N]$ and have
$\ell_N(p)=116N$, so their exponent is $116N-\kappa_N(p)$; the sum of $\kappa_N(p)\log p$ is
$N^2\int\kappa+o(N^2)$ by \eqref{eq:weighted-pnt}. The $2$-adic part contributes
$\mathsf A_2\log2\,N^2+o(N^2)$.
\end{proof}

The polynomial $\Delta_N$ has degree at most $n-g=d_0N$ and takes values in $D_N^{-1}\Z$ at the
integers. An integer-valued polynomial of degree $m$ has coefficients in $\frac1{m!}\Z$ (Newton's
interpolation formula), so
\begin{equation}\label{eq:QN}
 Q_N=(d_0N)!\,D_N\,\Delta_N(X)\in\Z[X],\qquad \deg Q_N\le d_0N,
\end{equation}
and $\log(d_0N)!=o(N^2)$. In this way the gain from the rank deficiencies enters through the integer
polynomial itself, and not through a separate common factor of its coefficients.

\begin{lemma}\label{lem:recycling}
For all large $N$ and every prime $N<q\le2N$, the $j$-th coefficient of $Q_N$ is divisible by $q^{2j}$.
\end{lemma}

\begin{proof}
For these primes the matrices $q^2\cA^0_N$ and $\cL_N$ are $q$-integral by Lemma~\ref{lem:digits}, and
$q^{2n}$ divides $D_N$. The identity
\[
 Q_N(X/q^2)=(d_0N)!\,\frac{D_N}{q^{2n}}\det\bigl(q^2\cA^0_N+X\cL_N\bigr)
\]
shows that $Q_N(X/q^2)$ has $q$-integral coefficients; its $j$-th coefficient is $c_{N,j}/q^{2j}$.
\end{proof}

\subsection{Nonvanishing}
\begin{proposition}\label{prop:nonvanishing}
Let $D$ be a design and suppose that the three $n_0\times n_0$ matrices $B_\sigma$, $\sigma\in\{0,1,-1\}$,
defined below, are invertible. Then for all large primes $p$ and every rational $z$ such that $pz$ has
denominator prime to $p$, $\Delta_p(z)\ne0$. In particular $\Delta_p(a/b)\ne0$ whenever $a\in\Z$, $b\ge1$
and $p^2\nmid b$.
\end{proposition}

The matrices are those of OpenAI's reduction modulo $p$ at $N=p$, with the design in place of $48$: let
$B$ be the rational $(n_0+1)\times n_0$ matrix of the entries of the filtered rational part at $N=1$ for
the rows $0,\dots,n_0$ (the row $n_0$ is an auxiliary row), and let $B_\sigma$ have rows
$B_r+\sigma B_{r+1}$, $0\le r<n_0$. For the design $D$, $B$ has $59$ rows and $58$ columns.

\begin{proof}
Let $p>4H_0$, with $H_0=H/N=78$, so that $H=H_0p<p^2/4$, and let $p$ be prime to the numerators and denominators
of the entries of $B$ and of the three determinants. In the construction at $N=p$ write the row index
$r=pr_0+i$ with $0\le r_0<n_0$ and $0\le i<p$. Since $q=q_0p$, $C=C_0p$, $h=h_0p$, $g=g_0p$, the Frobenius
relation $(1-t)^{p}\equiv1-t^p$ expresses the row polynomials modulo $p$ through those of the design at
$N=1$ and the successor row $r_0+1$; this is where the auxiliary row enters, and it is valid for rows on
both sides of $g$ because the polynomials are defined through $T_{r-g}$ and $U_{r-g-1}$ for signed
indices. The filter satisfies $s^{b_0p+k}(1-s)^{q_0p}=s^{\ell}(s^p)^{b_0+k_0}(1-s^p)^{q_0}$ in $\F_p[s]$, so
the filtered columns of one residue class use the raw columns $b_0+k_0,\dots,b_0+k_0+q_0$ of the design at
$N=1$ and never change the residue. By \cite[Lemma~3.1 and eqs.~(31), (35), (36)]{OpenAI2026}, which
need only $H<p^2/4$, the reduction of $p^2\cA^0_p$ modulo $p$ is then, after ordering the rows by $(i,r_0)$ and
the columns by $(\ell,k_0)$, the matrix $a^{\mathsf T}\otimes B_0+b^{\mathsf T}\otimes B_1$, where $a,b$ are the
fixed $p\times p$ matrices of OpenAI's identity (51), with $\det a=2^{-p(p-1)/2}$. The permutation
matrix $\Pi$ in this decomposition has the eigenvalue $0$ once and $\pm1$ each $(p-1)/2$ times, and
therefore
\[
 \det=(\det a)^{n_0}\,\det B_0\;\det(B_0+B_1)^{(p-1)/2}\det(B_0-B_1)^{(p-1)/2}\ne0
\]
in $\F_p$. Hence $p^{2n}\Delta_p(0)$ is a $p$-adic unit. For $z$ with $pz$ of denominator prime to $p$, the
difference $p^2z\cL_p=p\,(pz)\,\cL_p$ is $p$-integral with zero reduction (the entries of $\cL_p$ are
$p$-integral), so the reduction of $p^2(\cA^0_p+z\cL_p)$ is the same invertible matrix and
$v_p(\Delta_p(z))=-2n$. If $p^2\nmid b$, then $p(a/b)$ has denominator prime to $p$.
\end{proof}

For the design $D$ the invertibility of $B_0$, $B_0+B_1$ and $B_0-B_1$ is certified by exhibiting their
inverses modulo the prime $1009$, which exceeds $H_0=78$ so that every denominator of $B$ is a unit
there; the three determinants are then nonzero rational numbers. The computation modulo $1009$ certifies
the base matrices; the conclusion for all large primes is a theorem and not an extrapolation from
finitely many primes. Nothing in the proof uses the proportions of the design except $2g<n$, and $H<p^2/4$.

\section{Rank deficiencies modulo $p$}\label{sec:content}

In this section we prove Theorem~\ref{thm:content}. The principle is that of the arithmetic method of Hata
and Rhin--Viola, in the form that suits a determinant: if the reduction modulo $p$ of $p^2\cA^0_N$ has
rank at most $n-e$, then by Lemma~\ref{lem:rank-valuation} and Proposition~\ref{prop:valuation} the
determinant $\Delta_N(z)$ gains a factor $p^e$ at every $z$ with denominator prime to $p$, in
particular at every integer. The gain is in the denominator, and by \eqref{eq:QN} it is passed on to
the coefficients of $Q_N$. We therefore need lower bounds for the dimension of the left kernel of
$\cA^0_{N,p}$, that is, explicit row combinations $c\in\F_p^n$ with $c^{\mathsf T}\cA^0_{N,p}=0$. The
tools are the same throughout: the identities for Chebyshev polynomials in characteristic $p$
\begin{equation}\label{eq:cheb-p}
 \begin{gathered}
 T_p=X^p,\qquad U_{p-1}=(X^2-1)^{(p-1)/2},\\
 T_{kp}=T_k(X^p),\qquad U_{kp-1}=(X^2-1)^{(p-1)/2}\,U_{k-1}(X^p)
 \end{gathered}
\end{equation}
in $\F_p[X]$ (the last two follow from $T_{kp}=T_k\circ T_p$ and $U_{kp-1}=(U_{k-1}\circ T_p)\,U_{p-1}$), the relation
$(1-X)^p=1-X^p$, and the digit formula \eqref{eq:filtered-digits}.

\subsection{Setting}
Fix a design, an integer $N\ge1$ and an odd prime $p$ with $H<p^2$. For $0\le r<n$ put $m_r=r-g$, so that
$-g\le m_r\le n-1-g$, and $n-1-g\ge g$. In the notation of \eqref{eq:rows},
\[
 P_r=(1-X)^h\operatorname{rev}_{C-1}T_{m_r},\qquad D_r=(1-X)^h\operatorname{rev}_{C-1}U_{m_r-1},
\]
where $T_{-m}=T_m$ and $U_{-m-1}=-U_{m-1}$, so that $D_g=0$. For $c\in\F_p^n$ write
$S_c=\sum_rc_r\operatorname{rev}_{C-1}T_{m_r}$ and $U_c=\sum_rc_r\operatorname{rev}_{C-1}U_{m_r-1}$, so that $P_c=(1-X)^hS_c$
and $D_c=(1-X)^hU_c$. The lowest degree of $\operatorname{rev}_{C-1}T_m$ is at least $C-1-|m|\ge a:=A=b+q+2g$, and that of
$\operatorname{rev}_{C-1}U_{m-1}$ is at least $a+1$. Put
\[
 Q_c=(X-1)^q\,XWP_c,\qquad \widehat D_c=(X-1)^qD_c,\qquad W=(1-X^2)^{(p-1)/2},
\]
so that, with $\Lambda=g+q+h$, $Q_c$ has its coefficients in $[a+1,\,L+\Lambda+p-1]$ and $\widehat D_c$ in
$[a+1,\,L+\Lambda-1]$. Recall that a vector $c$ lies in the left kernel of $\cA^0_{N,p}$ as soon as the
right side of \eqref{eq:filtered-digits} vanishes for every filtered column $k$ (for the columns with a
carry among the $q+1$ raw columns one splits the sum over the raw columns at the carry; we shall not need
this case).

\subsection{Antisymmetric combinations of rows}
\begin{theorem}\label{thm:T}
Let $2g<p$, $H\le4p$ and $D_T=\min(g,\,n-1-g-2p,\,C-3p)\ge1$. Then
$\nu_N(p)\ge D_T-(n+q+g+h-3p)_+$.
\end{theorem}

The rows used are the differences of rows placed symmetrically about the three centres $g$, $g+p$ and
$g+2p$.

\begin{proof}
For $1\le d\le D_T$ and $k=0,1,2$ let $\delta_k(d)=e_{g+kp+d}-e_{g+kp-d}\in\F_p^n$. These rows exist because
$d\le g$ and $g+2p+d\le n-1$, and the $3D_T$ vectors $\delta_k(d)$ are linearly independent, since their supports
are disjoint: the supports lie in the intervals $[kp,kp+2g]$ of row indices, which are disjoint because
$2g<p$. The row $g+kp\pm d$ has $m=kp\pm d$. By $T_{a+b}-T_{a-b}=2(X^2-1)U_{a-1}U_{b-1}$ and
$U_{a+b-1}-U_{a-b-1}=2T_aU_{b-1}$ together with \eqref{eq:cheb-p},
\begin{align*}
 T_{kp+d}-T_{kp-d}&=2(X^2-1)^{(p+1)/2}\,U_{k-1}(X^p)\,U_{d-1},\\
 U_{kp+d-1}-U_{kp-d-1}&=2\,T_k(X^p)\,U_{d-1}.
\end{align*}
Put $R_d=2(1-X)^hX^{C-d}\operatorname{rev}_{d-1}U_{d-1}$, the value of $D_{\delta_0(d)}$. Reversing at degree $C-1$,
and using that $\operatorname{rev}_{kp}(T_k(X^p))=(\operatorname{rev}_kT_k)(X^p)$, one finds, with $T_0=1$, $T_1=X$ and
$T_2=2X^2-1$,
\[
 D_{\delta_1(d)}=X^{-p}R_d,\qquad D_{\delta_2(d)}=2X^{-2p}R_d-R_d,
\]
\begin{gather*}
 XWP_{\delta_0(d)}=0,\qquad XWP_{\delta_1(d)}=X^{-p}R_d-X^pR_d,\\
 XWP_{\delta_2(d)}=2X^{-2p}R_d-2R_d .
\end{gather*}
The polynomial $R_d$ is supported in degrees $[C-d,H-1]$, which lies in $[3p,4p)$ because $d\le C-3p$ and
$H\le4p$. So at each residue $\ell$ it has a single base-$p$ digit $\varrho_\ell(d)$, in position $u=3$.
For $x=(x_0,x_1,x_2)\in\F_p^3$ the response of the combination $\sum_kx_k\delta_k(d)$ at the raw column
$Jp+\ell$ is therefore, by Lemma~\ref{lem:digits}, $\varrho_\ell(d)\,f_J(x)$ with
\[
 f_J(x)=x_0\beta_{3,J}+x_1(\alpha_{2,J}-\alpha_{4,J}+\beta_{2,J})+x_2(2\alpha_{1,J}-2\alpha_{3,J}+2\beta_{1,J}-\beta_{3,J}),
\]
which depends neither on $d$ nor on $\ell$. The raw columns with $J\ge3$ do not meet the support, since
$3p+\ell\ge C-d\ge C-g=L$. Choose $x^*\ne0$ with $f_1(x^*)=f_2(x^*)=0$; this is two linear equations in three
unknowns. For $c=\sum_ds_d\sum_kx^*_k\delta_k(d)$ the responses at the columns with $J=1,2$ vanish, and the
columns with $J=0$, which have $\ell\ge b$, respond by $f_0(x^*)\sum_ds_d\varrho_\ell(d)$, the coefficient
of $X^{3p+\ell}$ in $\sum_ds_dR_d$. These are at most $(H-3p-b)_+=(n+q+g+h-3p)_+$ linear conditions on
$s=(s_d)$. The map $s\mapsto c$ is injective, so the solutions form a space of dimension at least
$D_T-(n+q+g+h-3p)_+$ in the left kernel.
\end{proof}

\subsection{A reflection of Chebyshev polynomials}
The first ingredient of the next theorem is a reflection formula modulo $p$.

\begin{lemma}\label{lem:u3-reflect}
Let $p$ be an odd prime and $\nu=(p+1)/2$. In $\F_p[X]$, for every integer $m$,
\[
 T_{p-m}=X^pT_m-(X^2-1)^{\nu}\,U_{m-1}.
\]
\end{lemma}

\begin{proof}
Both sides satisfy the three-term recurrence in $m$, so it suffices to check $m=0$ and $m=1$. This
needs \eqref{eq:cheb-p}. For $m=0$ it is $T_p=X^p$. For $m=1$ use $T_{p+1}=XT_p+(X^2-1)U_{p-1}$ and
$T_{p-1}=2XT_p-T_{p+1}$.
\end{proof}

\begin{theorem}\label{thm:U}
Let $2h\le p+1$, $L\le3p\le H$, $n+q+g+h\le3p$ and $2p\le n+q$, and put $n^*=\min(n,2p+2g+1)$. Then
$\nu_N(p)\ge n^*-g-2p-1$.
\end{theorem}

\begin{proof}
We restrict $c$ to the rows $r<n^*$, so that $m_r\le n^*-1-g\le2p+g$. Applying Lemma~\ref{lem:u3-reflect}
with $m=m_r$ and reflecting at degree $p+C-1$ gives, for every such row,
\[
 (1-X^2)^\nu\operatorname{rev}_{C-1}U_{m-1}-X\operatorname{rev}_{C-1}T_m=-X^{p+1}\operatorname{rev}_{C-1}T_{p-m},
\]
and the right side is divisible by $X^{p+C-|p-m|}$, hence by $X^L$ since $|p-m|\le p+g=p+C-L$. Multiplying
by $(1-X)^h$ and summing, for every $c$ supported on these rows
\begin{equation}\label{eq:identityK}
 X^L\mid(1-X^2)^\nu D_c-XP_c .
\end{equation}

Let $V$ be the space of such $c$ with $(1-X^2)^\nu\mid P_c$. Since $(1-X)^h$ divides every $P_r$, this is
$\nu-h$ conditions at $X=1$ and $\nu$ conditions at $X=-1$, that is, at most $p+1-h$ linear conditions.
For $c\in V$ write $P_c=(1-X^2)^\nu Y$ and $\widetilde Q=XY$. Since $(1-X^2)^p=1-X^{2p}$,
\[
 XWP_c=(1-X^{2p})\,\widetilde Q .
\]
The polynomial $P_c$ has no coefficients below $X^{A^*}$, $A^*=C-n^*+g$, and none from $X^H$ on, so
$\widetilde Q$ is supported in $[A^*+1,\,H-p-1]$; and dividing \eqref{eq:identityK} by the unit
$(1-X^2)^\nu$ modulo $X^L$ gives $D_c[k]=\widetilde Q[k]$ for $k<L$.

Now let $H-3p\le\ell<L-2p$. The inequality $L-2p\le A^*+1$ is equivalent to $n^*\le2p+2g+1$, and
$L-2p\le p$ since $L\le3p$. By the supports, $\widetilde Q[\ell]=0$, $\widetilde Q[2p+\ell]=0$ and
$\widetilde Q[3p+\ell]=0$, and $D_c[3p+\ell]=0$ because $3p+\ell\ge H$. So with
$\varkappa=\widetilde Q[p+\ell]$ the base-$p$ digits of $XWP_c$ at the residue $\ell$ are
$(0,\varkappa,0,-\varkappa,0,\dots)$ and those of $D_c$ are $(0,\varkappa,0,0,\dots)$. Hence, by
Lemma~\ref{lem:digits}, the response at every raw column $Jp+\ell$ with such an $\ell$ is $\varkappa$ times
a constant depending on $J$ only: on $V$ these columns, which have $J\in\{0,1,2\}$, span a space of
dimension at most one for each residue $\ell$.

Put $\ell_0=H-3p$ and $\ell_1=L-2p$; the hypotheses give $\ell_0\le b\le\ell_1$. For $\ell_0\le\ell<b$ the raw
columns $j\in[b,L)$ with this residue have $J\in\{1,2\}$, and for $b\le\ell<\ell_1$ they have
$J\in\{0,1,2\}$. So the saving is $2\ell_1-\ell_0-b$, and the number of linear conditions needed to make all
responses on the $n+q=L-b$ raw columns vanish on $V$ is
$(p+1-h)+(n+q)-(2\ell_1-\ell_0-b)$. Substituting $\ell_1=L-2p$, $\ell_0=H-3p$ this is $2p+1+g$. The space of
solutions in the span of $n^*$ rows has dimension at least $n^*-g-2p-1$.
\end{proof}

\subsection{A block of columns: the digits $1$ to $5$}
\begin{theorem}\label{thm:U6}
Let $e=(p+1)/2-q-h\ge0$, $n\le4p+2g+1$, $L\le6p$, and put $\mu_0=\max(q,H+q-6p)$, $\mu_1=L-5p$. Assume
$b+q\le p+\mu_0$. Then $\nu_N(p)\ge\mu_1-\mu_0-e$, that is
\[
 \nu_N(p)\ge\min\Bigl(\frac{p-1}2-g,\ n+b+q+h-\frac{11p+1}2\Bigr).
\]
The first term is the minimum if $6p\le H$ and the second if $6p\ge H$.
\end{theorem}

Unlike Theorems~\ref{thm:T} and~\ref{thm:U}, where the conditions are imposed on the rows, the argument
concerns a block of columns. We first need an identity.

\begin{lemma}\label{lem:K6}
In $\F_p[X]$, for every integer $k$, $2X^p(X^2-1)^{(p-1)/2}T_k-(2X^{2p}-1)U_{k-1}=U_{2p-k-1}$. If
$n\le4p+2g+1$, then for every $c$, $X^L\mid 2XWS_c-(2-X^{2p})U_c$, and consequently
\begin{equation}\label{eq:K6}
 2\,Q_c[i]=2\,\widehat D_c[i]-\widehat D_c[i-2p]\qquad(i<L).
\end{equation}
\end{lemma}

\begin{proof}
For all integers $m,k$ one has $U_{m-1}T_k-T_mU_{k-1}=U_{m-k-1}$, since both sides satisfy the same
recurrence in $k$ and agree for $k=0,1$. Take $m=2p$ and use \eqref{eq:cheb-p}:
$T_{2p}=2X^{2p}-1$ and $U_{2p-1}=2T_pU_{p-1}=2X^p(X^2-1)^{(p-1)/2}$. Reflecting at degree $2p+C-1$, the first
term becomes $2XW\operatorname{rev}_{C-1}T_k$, the second $(2-X^{2p})\operatorname{rev}_{C-1}U_{k-1}$, and the
right side $X^{2p+C-1}U_{2p-k-1}(1/X)$. For $k=m_r\ge-g$ this is divisible by $X^{C+k}$, hence by $X^L$, if
$2p-k-1\ge0$; and if $2p-k-1<0$ it has degree $k-2p-1$ and is divisible by $X^{4p+C-k}$, hence by $X^L$ when
$k\le4p+g$, which holds because $k\le n-1-g\le4p+g$. Summing with weights $c_r$ and multiplying by
$(X-1)^q(1-X)^h$ gives \eqref{eq:K6}.
\end{proof}

\begin{proof}[Proof of Theorem~\ref{thm:U6}]
Let $M=[\mu_0,\mu_1)$ and $w=\mu_1-\mu_0$; we may assume $w>e$. The block consists of the filtered columns
$k(J,\mu)=Jp+\mu-(b+q)$ for $1\le J\le5$ and $\mu\in M$. These indices are non-negative by the hypothesis
$b+q\le p+\mu_0$, and smaller than $n$ because $5p+\mu_1-1=L-1$; so there are $5w$ of them. As $\mu<p$, formula
\eqref{eq:filtered-digits} applies: the response of $c$ at $k(J,\mu)$ is
$\sum_u\alpha_{u,J}Q_c[up+\mu]+\beta_{u,J}\widehat D_c[up+\mu]$.

\emph{Supports.} For $\mu\in M$ we have $p+\mu\le p+\mu_1-1=L-4p-1\le a$ by $n\le4p+2g+1$, so the digits
$u=0,1$ of $Q_c$ and $\widehat D_c$ vanish. From $\mu_0\ge H+q-6p=L+\Lambda-6p$ we get that the digits $u\ge6$ of
$\widehat D_c$ and $u\ge7$ of $Q_c$ vanish. For $u\le5$ we have $up+\mu\le L-1$, so \eqref{eq:K6} applies.

\emph{From $Q$ to $\widehat D$.} For $u=2,\dots,5$, \eqref{eq:K6} gives
$\widehat D_c[up+\mu]=Q_c[up+\mu]+\tfrac12\widehat D_c[(u-2)p+\mu]$, and for $u=2,3$ the last index is at most $a$,
so the term vanishes. Hence if $Q_c[up+\mu]=0$ for $u=2,\dots,5$ and all $\mu\in M$, then by induction
$\widehat D_c[up+\mu]=0$ for $u=2,\dots,5$, and all digits of $\widehat D_c$ in $M$ vanish.

\emph{A relation from the Frobenius map.} We have
$Q_c=(-1)^qX(1-X)^{q+h+(p-1)/2}(1+X)^{(p-1)/2}S_c$, and $q+h+(p-1)/2+e=p$. Since $(1-X)^p=1-X^p$,
\[
 (1-X)^e\,Q_c=(1-X^p)\,Y_c,\qquad Y_c=(-1)^qX(1+X)^{(p-1)/2}S_c,
\]
and $Y_c$ is supported in $[a+1,\,C+(p-1)/2]$. Comparing the coefficients at $up+\mu$, subtracting the
coefficient at $(u-1)p+\mu$, and summing over $u=2,\dots,6$, the right side telescopes to
$Y_c[6p+\mu]-Y_c[p+\mu]$, which vanishes for $\mu\ge\mu_0+e$: the second term because $p+\mu\le a$, the
first because $6p+\mu\ge H+q+e=C+(p+1)/2>\deg Y_c$. So, if $Q_c[up+\mu']=0$ for $u=2,\dots,5$ and all
$\mu'\in M$, then
\[
 \sum_{t=0}^e(-1)^t\binom et\,Q_c[6p+\mu-t]=0\qquad(\mu_0+e\le\mu<\mu_1).
\]
This is a recurrence with unit leading coefficient: if $Q_c[6p+\mu]=0$ for the first $e$ values
$\mu=\mu_0,\dots,\mu_0+e-1$, then it vanishes for all $\mu\in M$.

\emph{Counting.} Impose on $c$ that the $n-5w$ filtered columns outside the block vanish, that
$Q_c[up+\mu]=0$ for $2\le u\le5$ and $\mu\in M$ ($4w$ conditions), and that $Q_c[6p+\mu]=0$ for the first $e$
values of $\mu$. By the three steps above all the responses of the block columns vanish. These are at
most $n-5w+4w+e=n-(w-e)$ linear conditions, so the left kernel has dimension at least $w-e$. Finally,
$w-e$ equals $\mu_1-q-e=n+b+q+h-(11p+1)/2$ if $\mu_0=q$ and $(p-1)/2-g$ if $\mu_0=H+q-6p$.
\end{proof}

\subsection{The second window}
\begin{theorem}\label{thm:S}
Let $e=(p+1)/2-q-h\ge0$ and $n\le4p+2g+1$, and put
\[
 m_0=\max(b+q,\,L+\Lambda-4p),\qquad m_1=\min(p,\,L-3p,\,a+1),\qquad a=b+q+2g .
\]
Then $\nu_N(p)\ge m_1-m_0-e$.
\end{theorem}

\begin{proof}
Let $M=[m_0,m_1)$ and $w=m_1-m_0>e$. For $\mu\in M$ we have $\mu\le a$, so the digit $u=0$ of $Q_c$ and
$\widehat D_c$ vanishes; $\mu\ge L+\Lambda-4p$ gives that the digits $u\ge4$ of $\widehat D_c$ and $u\ge5$ of
$Q_c$ vanish; and $3p+\mu<L$, so that \eqref{eq:K6} gives, with $D_i=\widehat D_c[ip+\mu]$ and
$Q_i=Q_c[ip+\mu]$,
\[
 Q_1=D_1,\qquad Q_2=D_2,\qquad Q_3=D_3-\tfrac12D_1 .
\]
So every response of the filtered matrix at a column with residue $\mu\in M$ is a combination of the four
functionals $D_1(\mu),D_2(\mu),D_3(\mu),Q_4(\mu)$ of $c$. For $\mu\in M$ the columns $k=\mu-(b+q)$,
$p+\mu-(b+q)$, $2p+\mu-(b+q)$ and $3p+\mu-(b+q)$ exist, since $b+q\le\mu<p$ and $3p+\mu<L$. So there are at least
$4w$ block columns, and their rows lie in the span $F$ of the $4w$ functionals.

We find $w-e$ linear relations among these functionals. As in the proof of Theorem~\ref{thm:U6},
$(1-X)^eQ_c=(1-X^p)Y_c$ with $Y_c$ a polynomial. Put $\sigma(\ell)=\sum_{u\ge0}Q_c[up+\ell]$ and sum the
coefficients of this identity over the residue class of $\mu\ge e$ modulo $p$. The right side gives
$\sum_u\bigl(Y_c[up+\mu]-Y_c[(u-1)p+\mu]\bigr)=0$, since $Y_c$ is a polynomial, and the left side gives
$\sum_{t=0}^e(-1)^t\binom et\sigma(\mu-t)$. For $\mu\in M$ with $\mu-e\ge m_0$ all $\mu-t$ lie in $M$, and there
$\sigma=Q_1+Q_2+Q_3+Q_4=D_1+D_2+D_3-\tfrac12D_1+Q_4$. So
\[
 \sum_{t=0}^e(-1)^t\binom et\bigl(\tfrac12D_1+D_2+D_3+Q_4\bigr)(\mu-t)=0\qquad(m_0+e\le\mu<m_1),
\]
and these $w-e$ relations are independent, since $Q_4(\mu)$ occurs with coefficient $1$ in the one for $\mu$
and only with lower arguments in the others. Hence $\dim F\le4w-(w-e)$, and the rank of the matrix is at most
$(n-4w)+3w+e=n-(w-e)$.
\end{proof}

The proof does not use a determinant or a lower bound for $N$. Unlike the proof of Theorem~\ref{thm:U6}, the
telescoping here is over all digits, so no vanishing of $Y_c$ at the ends is needed.

\subsection{The function $\kappa$ for the design $D$}
For $D$, $L=70N$, $H=78N$, $a=22N$, $\Lambda=11N$ and $q+h=6N$. Writing $x=p/N$, one checks the hypotheses of
the theorems on the intervals of the table of Section~\ref{sec:overview}. The hypotheses
$e\ge0$ and $n\le4p+2g+1$ both read $x\ge12$, and $L\le6p$ reads $x\ge35/3$; so Theorems~\ref{thm:U6}
and~\ref{thm:S} apply for $x\ge12$. For $12\le x\le13$ one has $\mu_0=H+q-6p=81N-6p$, so the first branch
of Theorem~\ref{thm:U6} applies, and for $x\ge13$ one has $\mu_0=3N$ and the second branch, which is positive
until $x=146/11$. In Theorem~\ref{thm:S}, $m_1=p$ for $x\le35/2$ and $m_1=70N-3p$ above, and
$m_0=81N-4p$ for $x\le69/4$ and $12N$ above; the value $m_1-m_0-e$ is $9x/2-75$, $x/2-6$ and
$64-7x/2$ on the three intervals $(50/3,69/4]$, $(69/4,35/2]$, $(35/2,128/7]$, and it is non-positive
outside $(50/3,128/7)$. Theorem~\ref{thm:T} applies for $x>10$ and $x\ge39/2$, and
$D_T=\min(5N,52N-1-2p,75N-3p)$; it gives $3p-(n+q+h)=3x-64$ up to $x=23$ and the cap $g=5N$ up to
$x=70/3$. From $x=70/3$ the hypotheses of Theorem~\ref{thm:U} hold up to $x=26$ (they read
$3p\ge L$, $3p\le H$, $2p\le n+q$ and $2h\le p+1$), and give $\min(n,2p+2g+1)-g-2p-1$, which is $5N$ for
$x\le24$ and $53N-2p-1$ above. In each case the value of the theorem is the value in the table minus at
most one unit, from the integer rounding. At a prime the largest of the values that apply is used.
The area of $\kappa$ is the sum of the areas of the eight cells, which are listed in
Section~\ref{sec:overview}.

\begin{remark}
Numerical computations of the exact rank of $\cA^0_{N,p}$ for several designs and primes show that the
bounds of Theorem~\ref{thm:content} are attained in the windows where they were tested, and that other
families add kernel where the windows overlap. These computations are not part of the proof. An upper bound
for the nullity is not proved, and none is needed.
\end{remark}

\section{The real place}\label{sec:real}

In this section we prove that $\Delta_N$ is small not only at $G$ but on an exponentially small interval
around it. The statement is Theorem~\ref{thm:real}: for all large $N$ and $|x-G|\le e^{-\rho N}$,
\[
 |\Delta_N(x)|\le\exp\bigl(-(n_0^2(U-\tfrac12\log2)-\varepsilon)N^2\bigr),\qquad \rho=n_0c_2+w .
\]
Here again the design is a parameter. We write $\alpha_A=A/n=(b_0+q_0+2g_0)/n_0$,
$\beta_h=2h_0/n_0$, $\gamma_H=H/n=(n_0+b_0+q_0+g_0+h_0)/n_0$, $\gamma_d=d_0/n_0=(n_0-g_0)/n_0$,
and
\[
 c_{07}=1-\frac{b_0+q_0+g_0}{n_0},\qquad c_{26}=3-\frac{b_0+q_0+2g_0}{n_0}.
\]
For $D$ these are $\alpha_A=11/29$, $\beta_h=3/29$, $\gamma_H=39/29$, $\gamma_d=53/58$, $c_{07}=41/58$ and
$c_{26}=76/29$. The proof has three parts. We first reduce the perturbation $\Delta_N(x)-\Delta_N(G)$ to
determinants in which some columns of $\cA_N$ are replaced by values of the row functions at fixed
points (\S\ref{ssec:quadrature}--\ref{ssec:andreief}). We then bound such determinants by OpenAI's
energy method, extended to configurations with fewer free points (\S\ref{ssec:energy}). Finally we
supply the trial sequences that the energy method needs and verify the resulting inequalities
(\S\ref{ssec:trials}--\ref{ssec:cells}).

\subsection{A quadrature formula for the coefficient of the constant}\label{ssec:quadrature}
Let $\delta=x-G$. Then $\Delta_N(x)=\det(\cA_N+\delta\cL_N)$, and the question is how large the
columns of $\cL_N$ are compared with those of $\cA_N$. By \eqref{eq:L} the entries of $\cL_N$ are linear
functionals of the coefficients of the row polynomials. Since these polynomials have bounded degree,
the functionals can be written exactly as finite combinations of point evaluations: if
$\tau_1,\ldots,\tau_J$ are distinct points and $J$ is at least the number of degrees of freedom of the
rows, Lagrange interpolation gives
\begin{equation}\label{eq:quad}
 \cL_N(r,k)=\sum_{a=1}^JA_r(\tau_a)\,\beta_{a,k},
\end{equation}
where the weights $\beta_{a,k}$ are the values of the same functionals on the Lagrange basis
polynomials. This is a Gauss-type quadrature formula, exact for polynomials of the relevant degree.
The point of the formula is that the columns of $\cL_N$ become combinations of point evaluations of the
same row functions that define $\cA_N$, so that the perturbed determinant can again be treated by the
energy method. The cost is the size of the weights.

We take $J=d_0N=53N$ nodes. Put $g_*=401/2500$ and $a_*=2g_*/(1+g_*^2)=2005000/6410801$, and let
$K=\lceil J/2\rceil$. The nodes are the points
\[
 t_{j,\pm}=\pm\sqrt{a_*^2+(1-a_*^2)(1+s_j)/2},\qquad T_K(s_j)=0,
\]
that is, the roots of $T_K\bigl((2t^2-(1+a_*^2))/(1-a_*^2)\bigr)$; we keep all $K$ positive roots and
the first $J-K$ negative ones, so that for odd $J$ exactly one root is dropped. They lie in
$(-1,-a_*)\cup(a_*,1)$. In the coordinate $x$ with $t=2x/(1+x^2)$ every node has $|x|>g_*$; this matters
below, because the energy fields are larger near $x=0$. Dropping a root changes the nodal polynomial by a
linear factor, and the derivative bound, the circle bound and the degree of the Chebyshev polynomial by
constant factors only, so an odd number of nodes costs nothing in the exponent. This is needed here,
since $d_0=53$ is odd and $J=53N$ is odd for odd $N$.

\begin{lemma}\label{lem:quadrature}
After removal of the common factor of the rows, the weights in \eqref{eq:quad} satisfy, for every
$w>\Wcap$ and all large $N$,
\[
 |\beta_{a,k}|\le108000N\,e^{wN}.
\]
\end{lemma}

\begin{proof}
At a root, the derivative identity for Chebyshev polynomials bounds $|P_J'(t_{j,\pm})|$ from below by a
constant multiple of $N a_*/(1-a_*^2)$, up to the dropped factor. The weight is a coefficient
functional of the Lagrange polynomial, and every coefficient of a polynomial is bounded by its maximum
modulus on a circle centred at the origin. On the circle $|z|=r$ with $r=1001/1000$, write $s=|\operatorname{Re}z|$.
For $\omega=(2z^2-(1+a_*^2))/(1-a_*^2)$ one has $|T_K(\omega)|\le(\Theta+\sqrt{\Theta^2-1})^K$, where
\[
 \Theta(s)=\frac{\sqrt{(r^2+1)^2-4s^2}+\sqrt{(r^2+a_*^2)^2-4a_*^2s^2}}{1-a_*^2}
\]
is half the sum of the distances from $\omega$ to $\pm1$. The function $\Theta$ decreases in $s$, while the
factor $|1-z|^{q_0+h_0}\le(r^2+1+2s)^{N(q_0+h_0)/2}$ coming from the column filter and from the common
factor of the rows increases. A cover of $[0,r]$ by $42$ intervals, with rational upper bounds for the
square roots at their endpoints, proves, with $d_0=53$ and $q_0+h_0=6$,
\[
 \frac{d_0}2\log\bigl(\Theta(s)+\sqrt{\Theta(s)^2-1}\bigr)+\frac{q_0+h_0}2\log(r^2+1+2s)<\Wcap\qquad(0\le s\le r),
\]
the supremum being $52.7849$ at $s\approx0.125$. The cover is done on the square of this inequality
raised to the power $N$, so that no logarithm is evaluated and the odd value of $d_0N$ costs only a
factor. Since $|z-t_{j,\pm}|\ge1/1000$ on the circle, the lemma follows.
\end{proof}

The circle of radius $1.001$ instead of $1$ separates the circle from the nodes; on the unit circle itself
the cancellation in the factor $(1-t)^h$ near $t=1$ is what keeps the bound small.

\subsection{Column expansion and an identity of Andr\'eief type}\label{ssec:andreief}
Expanding $\det(\cA_N+\delta\Phi\beta)$, where $\Phi(r,a)=A_r(\tau_a)$, multilinearly in the columns gives
a sum over maps $f$ from the columns to $\{\varnothing,1,\ldots,J\}$ of
$\prod_{f(k)\ne\varnothing}\delta\beta_{f(k),k}\det M_f$. Here $M_f$ is the matrix whose $k$-th column is the
$k$-th column of $\cA_N$ if $f(k)=\varnothing$ and the vector $(A_r(\tau_{f(k)}))_r$ otherwise.

\begin{lemma}\label{lem:expansion}
If $|\beta_{a,k}|\le B$ and $|\det M_f|\le Mc^{m(f)}$ for every $f$, where $m(f)$ is the number of columns
replaced, then
\[
 \bigl|\det(\cA_N+\delta\Phi\beta)-\det\cA_N\bigr|\le M\bigl((1+J|\delta|Bc)^{n}-1\bigr).
\]
\end{lemma}

Each $M_f$ has $p=n-m$ columns that are integrals and $m$ columns that are point values. Andr\'eief's
identity extends to this situation.

\begin{lemma}\label{lem:andreief}
Let $\mu,\nu$ be finite measures, $F_r$ $(r<p+m)$, $G_i$ $(i<p)$ and $K$ integrable functions, and let $B$
be the matrix with $B(r,i)=\iint F_r(t)K(t,s)G_i(s)\,d\mu(t)\,d\nu(s)$ for $i<p$ and
$B(r,p+j)=F_r(\tau'_j)$. Then
\begin{multline*}
 (p!)^2\det B=\int\!\!\int\det\bigl[F_r((t\Vert\tau')_c)\bigr]_{r,c}\\
 \cdot\det[K(t_i,s_j)]_{i,j<p}\det[G_i(s_j)]_{j,i<p}\,d\mu^p(t)\,d\nu^p(s),
\end{multline*}
where $t\Vert\tau'$ is the list of the $p$ variables followed by the $m$ fixed points.
\end{lemma}

\begin{proof}
For $m=0$ this is Andr\'eief's identity applied twice \cite{Forrester2018Andreief}. In general expand
$\det B$ multilinearly in its first $p$ columns, which introduces one pair of variables $(t_i,s_i)$ for
each of them, and antisymmetrize first in the $t$ and then in the $s$ variables; the last $m$ columns are
constant throughout.
\end{proof}

With $K(t,s)=1/(1-ts)$ the middle determinant is the Cauchy determinant, and the third is a
generalized Vandermonde determinant $\prod_js_j^b(1-s_j)^q\det[s_j^{E_i}]$. For the latter we use the
bound
\begin{equation}\label{eq:genvand}
 \bigl|\det[s_j^{E_i}]\bigr|\le\Bigl(\prod_j(1+s_j)\Bigr)^{n-p}\prod_{i<j}|s_i-s_j|,
\end{equation}
valid for distinct exponents $E_i<n$ and $s_j\in[0,1]$. It follows by multiplying the matrix by a
unitriangular matrix built from the polynomial $\prod_j(X-s_j)$ and bounding the remaining $(n-p)\times(n-p)$
determinant by the product of the $\ell^1$ norms of its columns. The quotient
$\det[s_j^{E_i}]/\prod_{i<j}(s_j-s_i)$ is a Schur polynomial, so the bound is close to optimal.

\subsection{The energy estimate with fixed points}\label{ssec:energy}
OpenAI bounds the integrand in Andr\'eief's identity by the exponential of $n^2E$, where $E$ is a discrete
logarithmic energy of the points $x_i$ and $s_j$ in the coordinate $t=2x/(1+x^2)$. The energy consists of
an external field for each point, logarithmic interactions $\log|x_i-x_j|$, $\log(1-x_ix_j)$ and
$\log|s_i-s_j|$ among points of the same kind, and an interaction $-\log(1-2x_is_j+x_i^2)$ between the two
kinds. There are two cases, according to which of OpenAI's two bounds for the determinant of the row
functions is used \cite[Props.~5.2, 5.3]{OpenAI2026}; in the second case the interaction among the $x_i$
has weight $2$ instead of $1$, and in the first the row field acquires a term with a multiplier
$\lambda_1$ for the constraint that defines the case.

The energy is bounded from above by duality. Given trial sequences $u=(u_k)$ and $v=(v_k)$, one writes
the interaction terms as quadratic forms in the moments $\frac1n\sum_iT_k(x_i)$ and
$\frac1n\sum_js_j^k$ and completes the squares against $u$ and $v$. What remains is the sum of the norms
$\|u\|^2=\sum_ku_k^2/k$, $\|v\|^2$ and, for each point, a function of one variable. With
$\mathsf T(u,x)=\sum_ku_kT_k(x)/k$ and $\mathsf S(v,x)=\sum_kv_kx^k/k$ these are, in the second case,
\begin{align*}
 X(x)&=\alpha_A\log|x|+\beta_h\log(1-x)-\gamma_H\log(1+x^2)-4\mathsf T(u,x)-\mathsf S(v,x),\\
 Y(s)&=\tfrac{b_0}{n_0}\log s+\tfrac{q_0}{n_0}\log(1-s)+2\mathsf T(v,s),
\end{align*}
and in the first case
\begin{multline*}
 X_1(x)=\alpha_A\log|x|+\beta_h\log(1-x)-(\gamma_H-\tfrac12)\log(1+x^2)\\
 +\lambda_1\bigl(\tfrac{2g_0}{n_0}-\tfrac{2x^2}{1+x^2}\bigr)
 -2\mathsf T(p,x)-\mathsf S(v,x),
\end{multline*}
with the same $Y$ (for the first case the trial sequence for the first variables is called $p$).
If $X\le M_X$ and $Y\le M_Y$ on the whole domains, the energy of any configuration is at most a
constant determined by $M_X$, $M_Y$ and the norms. In potential-theoretic terms the extremal
configuration is described by an equilibrium measure for two interacting measures in an external
field \cite{SaffTotik}, and the best trial sequences are its moments, as Carreira observed
\cite{Carreira2026}; any trial sequences give a valid upper bound.

When $m$ of the $n$ columns are point values, the configuration has only $n-m$ free points of each kind,
while the normalization is still by $n$. The fixed points contribute through the determinant of the row
functions evaluated at the concatenated list $t\Vert\tau'$. The fixed points are of two kinds: $a$ of
them are quadrature nodes, with $|x|>g_*$, and $b'$ of them are points $y\in(0,1)$; the numbers of free
points, of the first kind and of the second kind add up to $n$, and we put $\alpha=(a+b')/n$ and
$\tau=b'/n$. The missing free points change
the energy bound by an amount linear in $m$, which is the cost of a fixed column. Every step of OpenAI's
argument remains valid uniformly in the number of free points: the regularization of the logarithmic
interactions, the diagonal corrections, and the completion of squares only involve sums over the points
present. For the fixed points of the first kind the relevant field is the residual field
\[
 R(x)=-\gamma_d\log(1+x^2)-2\mathsf T(u,x),\qquad g_*\le|x|\le1,
\]
and for those of the second kind it is
\[
 N(y)=-\gamma_d\log(1+y^2)+\frac{n_0-2g_0-1}{n_0}\log y-u_1y,\qquad0<y<1,
\]
where $u_1$ is the first coefficient of $u$; the coefficient $(n_0-2g_0-1)/n_0$ is a lower bound, valid for
all $N\ge1$, of $(n-1-2g)/n$.

\begin{lemma}[budget, second case]\label{lem:budget2}
Let $\ell=\log2$, and let $\mathcal N,\Phi,\mathcal R,\mathcal N_c$ and $c_2$ be real numbers, with $m\ge0$,
such that
\[
 \mathcal N+\Phi+U-c_{07}\ell\le-m,\qquad(c_{26}\ell-\Phi+\mathcal R-c_2)^2\le8\ell m,\qquad
 \ell-\mathcal R+\mathcal N_c\le0 .
\]
Then for all $\alpha,\tau\ge0$
\begin{multline*}
 \bigl(-(c_{07}-\tfrac12)+(c_{26}-1)\alpha-\alpha^2+\tau-\tau^2\bigr)\ell+\mathcal N+(1-\alpha)\Phi+\alpha(1-\alpha)\ell\\
 +(\alpha-\tau)\mathcal R+\tau\mathcal N_c\le-U+\tfrac12\ell+c_2\alpha .
\end{multline*}
\end{lemma}

\begin{proof}
Subtracting the right side from the left gives
$(\mathcal N+\Phi+U-c_{07}\ell)+\alpha(c_{26}\ell-\Phi+\mathcal R-c_2)-2\ell\alpha^2
+\tau(\ell-\mathcal R+\mathcal N_c)-\ell\tau^2$. The last two terms are non-positive by the third hypothesis,
and the first three form a concave quadratic in $\alpha$ with non-positive discriminant, since
$A+s\alpha-2\ell\alpha^2\le0$ for all $\alpha$ as soon as $s^2+8\ell A\le0$.
\end{proof}

The lemma is applied with $\mathcal N$ a bound for $2\|u\|^2+\frac12\|v\|^2$, $\Phi=M_X+M_Y$, $\mathcal R$ a
bound for the supremum of $R$ and $\mathcal N_c$ a bound for the supremum of $N$. The analogue for the first
case has no fixed points of the second kind, and its polynomial in $\alpha$ is
\begin{multline*}
 \bigl(-(c_{07}-\tfrac12)+(c_{26}-1)\alpha-\alpha^2\bigr)\ell+n_1+(1-\alpha)F_1+\alpha(1-\alpha)\ell\\
 +\alpha\lambda_1\tfrac{2g_0}{n_0}\le-U+\tfrac12\ell+c_1\alpha,
\end{multline*}
with $n_1$ a bound for the norms of the first trial and $F_1=M_{X_1}+M_Y$; it holds for all real
$\alpha$ when its discriminant is non-positive and $c_1\le c_{26}\ell-F_1+\lambda_1\tfrac{2g_0}{n_0}$ (so that
it also holds for $\alpha<0$).

Combining the energy bounds with Lemma~\ref{lem:quadrature}, a determinant $M_f$ with $m$ fixed columns
contributes at most $(5\cdot108000N)^me^{(n_0c+w)Nm}$ times the unperturbed bound, where $c=\max(c_1,c_2)$;
the factor $5^m$ comes from expanding the row functions at the fixed points in the second case.
Lemma~\ref{lem:expansion} then gives, for $|\delta|\le e^{-\rho N}$ with $\rho=n_0c_2+w$, a factor
$(1+O(N^2))^{n}=e^{o(N^2)}$ (here $w$ is any number above $\Wcap$ and $c=c_2$). The unperturbed bound is
$\exp\bigl(n_0^2N^2(-U+\frac12\log2+\varepsilon)\bigr)$. This proves Theorem~\ref{thm:real}, given the
numerical bounds of the next two subsections.

\subsection{The trial sequences}\label{ssec:trials}
For the first case we use a trial sequence with $40$ Chebyshev coefficients and a finite number of
geometric terms, with the multiplier $\lambda_1=1.68493357$. For the second case we use trial sequences
of the form
\begin{equation}\label{eq:trial}
 u_k=c_k+\operatorname{Re}\sum_jr_jz_j^k\qquad(k\ge1;\ c_k=0\text{ for }k>K),
\end{equation}
with $K=\Kcoef$ Chebyshev coefficients, dyadic rationals, and $\Ngeo$ geometric terms, $\Ncomplex$ of them with
complex $z_j$ and complex $r_j$.

The geometric terms have a simple meaning. Since $\sum_{k\ge1}z^kT_k(x)/k=-\frac12\log(1-2xz+z^2)$, a term
$(z,r)$ contributes $-\frac12\operatorname{Re}\,r\log(1-2xz+z^2)$ to the potential $\mathsf T(u,x)$. If
$x_0+i\epsilon=\frac12(z+z^{-1})$ then $1-2xz+z^2=2z(x_0+i\epsilon-x)$, so the term is, up to a constant, the
logarithmic potential of a point charge of complex strength $r$ at $x_0+i\epsilon$ and its conjugate. The
extremal measure of the energy problem has a gap in its support around $x=0$, and its potential has a
logarithmic singularity at $x=0$ and square-root behaviour at the edges of the gap. A Chebyshev expansion
converges slowly near such points. A charge at a small distance $\epsilon$ from them reproduces the local
behaviour with one term. We place the complex charges at
$x_0\in\{0,\pm0.03,\pm0.06,\pm0.09,\pm0.12,\pm0.16\}$ with $\epsilon\in\{10^{-2},3\cdot10^{-3}\}$, and
the real terms at points $\pm z$ with $z$ increasing from $0.8$ to $0.999997$; the numbers $z_j$ are
rationals with at most twelve decimals. The norms of the geometric terms have closed forms: two terms $(z_j,r_j)$,
$(z_l,r_l)$ contribute
$-\frac12\operatorname{Re}\bigl(r_jr_l\log(1-z_jz_l)+r_j\bar r_l\log(1-z_j\bar z_l)\bigr)$ to $\|u\|^2$.

\subsection{Verification of the bounds}\label{ssec:cells}
We must bound three functions of one variable from above on their whole domains: the field $X$ on
$[-1,1]$, the field $Y$ on $[0,1]$, and the residual field $R$ on $g_*\le|x|\le1$. With $x=\cos\theta$,
each is a constant plus a cosine series with $K$ terms plus finitely many terms
$\operatorname{Re}\,w\log(1-ae^{im\theta})$.

We cover each domain by cells $[x_c-h,x_c+h]$ whose endpoints are multiples of $2^{-40}$, and on
each cell bound the function by its Taylor expansion at $x_c$. The cells are grouped in blocks. The
polynomial part has degree $K-1$, and a global bound for its fifth derivative would force very short cells.
Instead, for each block the derivatives of the polynomial part are computed exactly at the centre of the
block up to order $70$, and the remainder beyond order $70$ is bounded by a single constant for the block,
which enters multiplied by $h^{70}/70!$; the Taylor expansion at the centre of a cell is obtained from
that of the block by an exact shift. The logarithmic terms are expanded to order $5$, with remainders
bounded from the exact minimum of $|1-ae^{im\theta}|$ on the cell. A few cells at $x=0$ and $x=\pm1$ treat
the singular terms separately. The covers of $X$, $Y$ and $R$ have \Nblocks{} blocks and \Ncells{} cells, respectively.
Each cell is checked by exact rational and integer arithmetic.

The resulting bounds are
\[
 \sup X\le\SupX,\quad \sup Y\le\SupY,\quad
 \sup_{g_*\le|x|\le1}R\le\ResCap,
\]
so that $M_X+M_Y=\FieldCap$; the norms satisfy $2\|u\|^2+\frac12\|v\|^2\le\NormCap$; and the first
coefficient satisfies $u_1\ge9/50$, which is used near $x=0$. For the near points, a one-variable
estimate shows that
$\sup_{0<y<1}N(y)\le-77/100$ for $u_1\ge9/50$, with the weight $47/58$ in $N(y)$; the true supremum is
about $-0.7758$. For the first case, with the weights $(\alpha_A,\beta_h,\gamma_H-\frac12,\lambda_1,2g_0/n_0)=(11/29,3/29,49/58,\lambda_1,5/29)$,
the bounds $\sup X_1\le-1.15854$ and $\sup Y\le-1.32981$ are obtained by isolating the zeros of the
numerators of the derivatives, which are polynomials with integer coefficients, and bounding the values
between them; the norm bound is $n_1=0.67761$.

With $c_{07}=41/58$, $c_{26}=76/29$ and $U=\Utarget$, the hypotheses of Lemma~\ref{lem:budget2} hold with
\begin{gather*}
 m=\Mmargin,\qquad \ell-\mathcal R+\mathcal N_c\le-0.0137,\\
 (c_{26}\ell-\Phi+\mathcal R-c_2)^2\le\Disc\le8\ell m=\DiscBound,
\end{gather*}
with $c_2=\Ctwo$; for the first case $c_1=\Cone<c_2$ is admissible. This gives the cost $c_2$ of a fixed
column, and $\rho=58c_2+w=\Rho$ for $w=2.9\cdot18.203=52.7887$ ($>\Wcap$). The decay rate in
Theorem~\ref{thm:real} is $n_0^2(U-\frac12\log2)=\TDelta\ldots$ The trial data are fixed once and for
all; the file containing them has SHA-256 digest
\begin{center}
{\footnotesize\texttt{\Digest}}.
\end{center}

\section{Proof of Theorem~\ref{thm:main}}\label{sec:proof}

\subsection{The criterion}
We prove Proposition~\ref{prop:criterion}. Its proof shows where each of the three properties of
$Q_N$ enters: the nonvanishing and the smallness are used once at the end, and the divisibility by prime
squares is what makes the choice of the prime free.

\begin{proof}[Proof of Proposition~\ref{prop:criterion}]
Let $a/b$ be a rational number with $b$ large, put $L=\log b$, $U=\kappa L$ and $V=(\kappa-\omega)L$, and
suppose $|\xi-a/b|\le b^{-\rho\kappa}$. By the prime number theorem, $(1-\delta)t\le\theta(t)\le(1+\delta)t$
for all large $t$, and $V\to\infty$ with $b$. If every prime $p\in(V,U]$ had $p^2\mid b$, then
$2(\theta(U)-\theta(V))\le\log b=L$; but
\[
 2\bigl(\theta(U)-\theta(V)\bigr)\ge2L\bigl((1-\delta)\kappa-(1+\delta)(\kappa-\omega)\bigr)
 =2L\bigl(\omega-\delta(2\kappa-\omega)\bigr)>L ,
\]
by the fourth hypothesis. So there are primes in $(V,U]$ whose squares do not divide $b$. Let $p$ be the
\emph{largest} one. Every prime $q\in(p,U]$ has $q^2\mid b$, and $q\le U<2V<2p$ because $2\omega<\kappa$.
Let $B$ be the product of these $q^2$; then $B\mid b$ and $\log B=2(\theta(U)-\theta(p))$.

Every term $c_ia^ib^{dp-i}$ of the integer $b^{dp}Q_p(a/b)=\sum_ic_ia^ib^{dp-i}$ is divisible by $q^{2dp}$ for each
of these $q$: $q^{2i}$ divides $c_i$ by hypothesis, and $q^2\mid b$ gives $q^{2(dp-i)}\mid b^{dp-i}$. So $B^{dp}$
divides the integer. Since $Q_p(a/b)\ne0$ (as $p^2\nmid b$), $b^{dp}Q_p(a/b)/B^{dp}$ is a nonzero integer, and
\begin{equation}\label{eq:recycled}
 1\le\Bigl(\frac bB\Bigr)^{dp}\,|Q_p(a/b)| .
\end{equation}
Since $\rho p\le\rho U=\rho\kappa L$, we have $|a/b-\xi|\le e^{-\rho p}$, so $|Q_p(a/b)|\le e^{-\eta p^2}$ and
\eqref{eq:recycled} gives $\eta p^2\le dp\bigl(L-2(\theta(U)-\theta(p))\bigr)$, that is,
\[
 \eta p\le d\bigl(L-2(\theta(U)-\theta(p))\bigr)\le d\bigl(L-2U+2p+4\delta U\bigr),
\]
using $\theta(U)\ge(1-\delta)U$ and $\theta(p)\le(1+\delta)p\le(1+\delta)U$ in the form
$\theta(U)-\theta(p)\ge U-p-2\delta U$. Hence $(\eta-2d)p\le dL-2dU+4d\delta U$. As $\eta<2d$ and $p\le U$,
the left side is at least $(\eta-2d)U$, so $\eta U\le dL+4d\delta U$, that is, $\eta\kappa\le d+4d\delta\kappa$,
which contradicts the last hypothesis.
\end{proof}

The choice of the largest admissible prime is what makes the selection free: the primes of the window
that are not admissible have their squares in $b$, and these squares are recovered from the coefficients
through the divisibility hypothesis. A cruder choice, of any prime in a short interval that does not
divide $b$, loses a factor. The free lower end $V=(\kappa-\omega)L$ is what allows $\kappa<2$: the condition $q\le U<2p$ for the recycled
primes needs only $U<2V$.

\subsection{The hypotheses for the design $D$}
Take the design $D=(58,9,3,5,3)$ and $d=53$ (the degree of $Q_N$ is at most $53N$, and the polynomials
$Q_p$ of the proposition are the $Q_N$ with $N=p$), and
\begin{gather*}
 \rho=\Rho,\quad \eta=\EtaTot,\quad \omega=\tfrac34,\quad \delta=2.9\cdot10^{-7},\\
 \kappa=\frac{\Bbank}{\rho}=\Kappa .
\end{gather*}
We check the hypotheses.
\begin{itemize}
\item $\deg Q_N\le53N$ by \eqref{eq:QN} and Lemma~\ref{lem:degree}.
\item The nonvanishing: by Proposition~\ref{prop:nonvanishing}, $\Delta_p(a/b)\ne0$ for large $p$ with
$p^2\nmid b$, hence $Q_p(a/b)\ne0$.
\item The divisibility: by Lemma~\ref{lem:recycling}, $q^{2i}$ divides the $i$-th coefficient of $Q_p$ for all
primes $p<q\le2p$.
\item The smallness: by Theorem~\ref{thm:real}, Proposition~\ref{prop:denominators} and $\log(53N)!=o(N^2)$,
for $|x-G|\le e^{-\rho N}$ and large $N$,
\[
 |Q_N(x)|\le\exp\bigl(-(\TDelta\ldots-6395.7087\ldots-3\varepsilon)N^2\bigr)\le e^{-\eta N^2}
\]
since $\eta=\EtaTot<\EtaExact\ldots$ is rounded down.
\item The numerical conditions: $\rho\kappa=\Bbank$; $0<\eta<2d=106$; $\omega=3/4\le1$ and $2\omega=1.5<\kappa=\Kappa$;
$2(\omega-\delta(2\kappa-\omega))=1.49999\ldots>1$; and
$d+4d\delta\kappa=\PayL\ldots<\eta\kappa=\PayR\ldots$.
\end{itemize}
Proposition~\ref{prop:criterion} therefore gives $|G-a/b|>b^{-\rho\kappa}=b^{-\Bbank}$ for all integers
$a$ and all sufficiently large $b$. This proves Theorem~\ref{thm:main}. The last condition is
equivalent to $\Bbank>x=d\rho/(\eta-4d\delta)=\Bx$, and in the limit $\delta\to0$ the argument gives
the exponent $d\rho/\eta=\Blimit\ldots$ for the constants above; the integer $\Bbank$ is the one for
which all conditions are verified with rational constants.

\section{Remarks on the method}\label{sec:remarks}

The exponent $d\rho/\eta$ has three ingredients. The degree $d=n-g$ and the perturbation exponent
$\rho=n_0c_2+w$ come from the construction and from the quadrature: of $\rho=\Rho$, the cost of fixed
columns $n_0c_2$ is $\NCtwo$ and the logarithmic size of the weights is $52.79$. The decay rate $\eta$ is
a small difference of two large numbers, the energy at the real place, $\TDelta$, and the exponent of the
denominators, $6395.71$, and the rank deficiencies of Theorem~\ref{thm:content} contribute $19.14$ of the
$\EtaTot$. Since the exponent is inversely proportional to $\eta$, an additional gain of $0.01$ in the area
of $\kappa$ lowers it by about $0.22$.

All four theorems of Section~\ref{sec:content} are lower bounds for the nullity, and the proofs give the
kernel vectors explicitly. They are proved for every design, and so are the other arguments of the paper; a
different design changes $d$, $\rho$ and $\eta$ together, and we make no claim that $D$ is optimal. In
particular the bound $\kappa$ of Theorem~\ref{thm:content} does not use any relation between different
families, and at a prime where windows overlap the exact nullity is larger than each bound. An upper bound
for the nullity, or a proof that the kernels of two families are independent, would increase $\eta$ without
changing anything else.

\section{Formalization}\label{sec:formal}

All results of this paper are formalized in Lean~4 \cite{Lean4} with Mathlib \cite{Mathlib}, on top of
OpenAI's formalization of the irrationality of $G$ \cite{OpenAI2026}. The development is available at
\begin{center}
\FormalizationLink.
\end{center}
The statement of Theorem~\ref{thm:main} in \lean{Challenge.lean} uses only Mathlib, with Catalan's constant
written as the series $\sum_j(-1)^j/(2j+1)^2$; it is the theorem
\begin{center}
\texttt{CatalanIrrationalityMeasure.catalan\_irrationalityMeasure\_\Bbank}
\end{center}
and its proof is \lean{Solution.lean}. The solution has no \lean{sorry} and no \lean{native\_decide}, and the
command \lean{\#print axioms} reports only the three standard axioms \lean{propext}, \lean{Classical.choice} and
\lean{Quot.sound}. The file \lean{comparator.json} is a configuration for \lean{lake comparator}
\cite{Comparator} stating that \lean{Solution} proves the statement of \lean{Challenge} with these axioms
only. The prime number theorem is taken from the PNT+ project \cite{PNTPlus}. The parts of OpenAI's
formalization that are used were converted to Lean's module system and ported to a recent Lean release; the
repository's \texttt{README} and \texttt{PORTING.md} describe the changes.

\subsection{Module layout}
Everything specific to this paper is in the directory \lean{GenD} of \lean{CatalanMeasure}. Its design is that of the paper:
a design is data, the generic modules prove every statement for an arbitrary design, and the instance
files contain numerals only.
\begin{description}[leftmargin=1.5em,style=nextline]
\item[Generic layer] \lean{PL} (piecewise linear profiles, their areas) and \lean{PrimeGrowth} (the
weighted prime number theorem \eqref{eq:weighted-pnt} for ramp profiles; one exponent per prime);
\lean{Budget} (the scalar inequalities of Lemma~\ref{lem:budget2} and its analogue for the first case, and the
conditions of the criterion), \lean{PayW} and \lean{Recycled} (Proposition~\ref{prop:criterion} with the free
window end); \lean{Final} (the structure \lean{DesignData} and the theorem \lean{measure\_of\_design}, which
contain no numeral).
\item[Denominators] \lean{OddProfile} (the valuations $\ell_N(p)$ for the lower and upper pools),
\lean{TwoAdic} (the bound \eqref{eq:A2}), \lean{DenomExp}, \lean{DenomGrowth}, \lean{DenomGrowth2},
\lean{Denominator} and \lean{Normalizer} (the common denominator, $Q_N$, the divisibility of
Lemma~\ref{lem:recycling}).
\item[Rank deficiencies] \lean{Interface} (the nullity, the four families and \lean{famKap\_le\_nullity}),
\lean{RankCommon}, \lean{RankK6}, \lean{RankDigits}, \lean{RankFam}, \lean{RankFilt}, \lean{RankSeq},
\lean{RankBlock} (the shared layers: Chebyshev reflection, identity \eqref{eq:K6}, the digit formula in
filtered form, the block framework), and \lean{RankT}, \lean{RankU}, \lean{RankUCore}, \lean{RankU6},
\lean{RankS} for Theorems~\ref{thm:T}--\ref{thm:S}; \lean{KapHelpers} turns the comparison with a window list
into linear arithmetic.
\item[Nonvanishing] \lean{Nonvanishing} (Proposition~\ref{prop:nonvanishing} from the three base matrices),
and the residue-table certificates \lean{CertMod}, \lean{ResBase}, \lean{ResCheck}, \lean{ResInv},
\lean{ResRec}.
\item[Real place] \lean{Analytic/} (28 modules): the constants and fields of a design (\lean{Const},
\lean{Fields}, \lean{CommonField}), the first-case and second-case contracts (\lean{First},
\lean{FirstGlue}, \lean{Second}, \lean{Glue}), the grouped products over the fixed points (\lean{GroupedProduct},
\lean{MarkedSheet}, \lean{MarkedSheetOn}, \lean{MarkedMixedMajorant}, \lean{MarkedSecond}), the prefactor,
the perturbation by the quadrature nodes (\lean{Prefactor}, \lean{AtomBase}, \lean{WeightedAtom},
\lean{DeltaSmall}, \lean{DeltaChain}), and the nodes and the weight bounds for any number of nodes
(\lean{GapNodes}, \lean{GapNodal}, \lean{GapWeights}, \lean{FullWeights}, \lean{GapCircleGen},
\lean{CircleCover}). \lean{Cert/} (6 modules) is the generic cell machinery: the evaluator of logarithms
(\lean{WLog}, \lean{WLogSpec}), the residual field with a weight (\lean{ResW}), the shifted tables
(\lean{PolyShift}) and the assembly of cells and blocks (\lean{CellLoc}, \lean{Assembly}).
\item[The design $D$] \lean{D58/} holds the instance: \lean{Budget} (the constants of the certificate),
\lean{Ramps}, \lean{Kap}, \lean{OddProfile}, \lean{TwoAdic}, \lean{Normalizer}, \lean{Pay}, \lean{Final}
(the design, the window list, the profile $f$, $\mathsf A_2=2673/8$, the growth of the denominator and
the conditions of the criterion), \lean{Circle} (the 42-piece cover), \lean{FirstTrial},
\lean{FirstTrialCert}, \lean{FirstBridge} and \lean{FirstCert/} (the first case), \lean{Cert/} (the second
case: trial data, the covers of $X$, $Y$ and $R$ in blocks, the norms, the near-field cap),
\lean{SecondTrial}, \lean{C3/} (the three inverses modulo $1009$), \lean{Certified} and \lean{Main}, which
proves \lean{catalan\_irrationalityMeasure\_\Bbank} in the namespace \lean{CatalanMeasure.GenD.D58}.
\end{description}
The public theorem in \lean{Solution.lean} is a one-line wrapper of that theorem. The modules of
OpenAI's formalization and the generic valuation theorems of the construction
(\lean{CatalanMeasure/Parametric}) are used as they are.

\subsection{Certificates}
The proof combines conventional formal proofs with numerical certificates checked by the Lean kernel:
the inverses of the base matrices modulo $1009$, the cover of Lemma~\ref{lem:quadrature}, the trial
sequences of Section~\ref{ssec:trials} and the covers of Section~\ref{ssec:cells}. All are exact rational or
integer computations, decided by the kernel without compiled code. The interval arithmetic used to
generate the trial data and the covers \cite{Arb} is not trusted: Lean rechecks every inequality. Four
devices keep the certificates small enough for the kernel.
\begin{itemize}
\item \emph{Shared Taylor bounds.} The cells of a cover are grouped in blocks. For each block the
derivatives of the polynomial part of the field are evaluated once, exactly, at the centre of the block,
together with one bound for the remainder of the expansion; every cell of the block obtains its expansion
by an exact shift, instead of repeating the evaluation of a polynomial with $\Kcoef$ coefficients.
\item \emph{Accuracy per term from a shared table of nodes.} The logarithms are evaluated by the series
$\log\frac{1+x}{1-x}=2x\sum_i x^{2i}/(2i+1)$ with an explicit geometric tail, and the number of terms is
chosen for each logarithmic term separately from a table of nodes shared by all the cells, instead of
one accuracy for the whole sum.
\item \emph{Residue recurrences.} The tables of the moments $M(i,j)$ and $Z(i,j)$ modulo $1009$ are not
computed from the rational numbers. They are generated by the one-step recurrences that define these
numbers, once proved correct, and the kernel checks the recurrence, the filtered rows against the
row polynomials and the products of the matrices with their claimed inverses, row by row, in natural
numbers.
\item \emph{$2$-adic normalisation.} Where a recurrence produces integers divisible by large powers of
$2$, as the Chebyshev and Clenshaw recurrences do, the tables store the quotients by these powers, with
one lemma relating the two; this also avoids integers that the kernel's hashing of literals treats as
equal.
\end{itemize}

\section*{Disclosure of AI use}

The mathematics, the Lean formalization, the numerical computations and certificate generators, the
literature search and the text of this paper were produced by AI systems at the request of the author, who
posed the problem, directed the work, and is responsible for this paper. By component: Claude
(Opus~5.5 and Sonnet~5.5, Anthropic) produced the treatment of an arbitrary design (the degree, the
denominators at all primes and the nonvanishing, Section~\ref{sec:setup}), the digit formula and the rank
deficiencies of Theorems~\ref{thm:T}, \ref{thm:U}, \ref{thm:U6} and~\ref{thm:S}, the passage from the
nullity to the denominators, the criterion with a free lower end of the window, the perturbation estimate
by quadrature with an odd number of nodes, the trial sequences with complex terms and the verification of
Section~\ref{ssec:cells}, the choice of the design and the final assembly, together with the corresponding
Lean proofs and the generic layer of Section~\ref{sec:formal}. OpenAI Codex produced the energy estimates with
fixed columns and the fields of the first case, the identity of Andr\'eief type with fixed points, the
bound for the weights of Lemma~\ref{lem:quadrature}, the divisibility of the coefficients by prime squares
and the choice of the largest admissible prime, with the corresponding Lean proofs for one particular
choice of the parameters, which Claude generalised. The underlying irrationality proof and its
formalization are OpenAI's \cite{OpenAI2026}.

\subsection*{Acknowledgements}
We thank the authors of Mathlib and of the PNT+ project, and OpenAI for making its formalization available.


\bibliographystyle{amsplain}
\bibliography{refs}

\end{document}
